\chapter{图像生成与编辑}
\label{chap:vision-image-generation}

设想一位设计者手里有一只红色玩具狗的几张照片，提出四个请求：“画这只玩具狗坐在蓝色杯子左边，在雪地里”；“把这张照片转成水彩画”；“把照片里狗旁边的杯子换成花瓶”；“根据绕着它拍的一圈照片，给出从正后方看的样子”。四个请求的输出都是一幅新图像，但它们给定的东西不同，允许变化的东西也不同。第一个请求只规定内容和主体，构图、光照和背景细节都可以由生成模型补充；第二个请求有一幅输入，狗的姿态与数量要保持，笔触和色调可以整体改变；第三个请求同样有输入，却只允许杯子所在的区域改变；第四个请求的答案受真实拍摄约束，背面若没有拍到，模型补出的内容就不能算作观测结果。

本章要回答的问题是：给定这些不同的条件，怎样得到满足要求的图像，又怎样核验它确实满足了要求。生成模型的训练与采样原理已在前面各卷建立：扩散模型通过学习逐步去噪，从高斯噪声出发采样出图像，可参考\BookChapterRef{chap:pgm-generative}{卷六的“概率生成模型”章}；生成对抗网络以判别器的反馈训练生成器，条件生成对抗网络和配对、无配对的图像转换见\BookChapterRef{chap:neural-generative-models}{卷五的“基于神经网络的生成模型”章}。本卷\ref{chap:vision-multimodal-models}章又讲了LDM、unCLIP、GLIDE与流匹配等图像生成模型的结构和训练。本章不再重讲这些模型，而是研究它们在任务中怎样使用，并完整讲解只服务于某一任务的模型，例如空间控制用的ControlNet、主体适配用的DreamBooth和新视角合成用的NeRF。

四个请求依次放宽“可以变化”的范围，也依次加强“必须保持”的依据，这决定了本章的路线。我们先把四类任务的输入、输出和评价指标写清楚，并从这四个请求中归纳出贯穿全章的四个挑战：一个条件对应许多合格图像，多个对象的属性容易互相串换，局部修改与内容保持彼此牵制，有限的照片难以保证任意视角下的一致。新图像的创建只受条件约束，难点是把文字、空间图和主体参考送进生成过程，并从多个候选中选出满足要求的一幅；已有图像的转换多了一幅输入，难点变成怎样学到保持对应关系的映射；指定内容的编辑进一步要求区域外不变，这里先讲泊松图像编辑、基于样例的修补和PatchMatch等经典做法，它们只从已有图像中搬运内容，再讲按文字生成新内容的学习方法，还要在多轮请求之间维护已接受的状态；新视角合成则不再依靠生成先验，而是在多视图几何的基础上，从多幅照片拟合一个可以从任意相机渲染的场景。最后一节把这四类任务的核验放在一起，按条件遵循、输入保持、观测支持与成本三个维度比较。贯穿全章的玩具狗、杯子及各类分数和坐标都是教学构造，不代表任何模型的实测结果。

\section{图像生成与编辑的问题与挑战}
\label{sec:vision-generation-challenges}

开篇四个请求的输出都是图像，输入却各不相同。新图像的创建以文字、空间图或主体照片为条件，输出一幅或几幅满足条件的图像；已有图像的转换以一幅图像为输入，输出同一内容的另一种呈现形式；指定内容的编辑以原图、编辑请求和可能附带的范围标注为输入，输出只在请求涉及的部分发生变化的图像；新视角合成以多幅带相机参数的照片为输入，输出指定新相机下的图像。前三类任务的合格答案通常不止一个，新视角合成原则上只有一个答案，即在那个相机位置真实拍到的照片。这一差别决定了评价方式：前三类任务只能检查结果是否落在可接受的范围里，新视角合成可以逐像素与真实照片比较。

评价指标按所核验的对象分为四组。第一组比较生成图像整体的分布与真实图像的分布。\term{Fréchet Inception距离}（Fréchet Inception distance, FID）把一组真实图像和一组生成图像分别送入在ImageNet上训练的Inception-v3网络，取最后一个池化层的2048维特征，为两组特征各拟合一个高斯分布，再计算两个高斯之间的Fréchet距离\scite{vision-fid2017}：
\begin{equation}
 \mathrm{FID}
 =\lVert\bm{\mu}_{\mathrm{r}}-\bm{\mu}_{\mathrm{g}}\rVert_2^2
 +\operatorname{tr}\!\Bigl(
  \bm{\Sigma}_{\mathrm{r}}+\bm{\Sigma}_{\mathrm{g}}
  -2\bigl(\bm{\Sigma}_{\mathrm{r}}\bm{\Sigma}_{\mathrm{g}}\bigr)^{1/2}
 \Bigr)\text{，}
 \label{eq:vision-fid}
\end{equation}
其中$\bm{\mu}_{\mathrm{r}},\bm{\Sigma}_{\mathrm{r}}$和$\bm{\mu}_{\mathrm{g}},\bm{\Sigma}_{\mathrm{g}}$分别是真实特征与生成特征的均值和协方差，$(\cdot)^{1/2}$表示矩阵平方根。第一项比较两组特征的中心，第二项比较它们的分散程度和各维之间的相关；FID越小，两组特征的分布越接近。一维时公式化简为$(\mu_{\mathrm{r}}-\mu_{\mathrm{g}})^2+(\sigma_{\mathrm{r}}-\sigma_{\mathrm{g}})^2$。若真实特征的均值为0、标准差为1，生成特征的均值为0.5、标准差为2，则FID为$0.25+1=1.25$；只把生成特征的标准差调回1，FID降为0.25，说明这个指标会同时惩罚位置偏移和多样性的偏差。这里的Fréchet指标是两个高斯之间2阶Wasserstein距离的平方；一般协方差下的闭式及其一维化简可参考\BookChapterRef{chap:information-theory-and-statistical-geometry}{卷2《数学准备：概率、因果与决策》的“概率分布的分析与优化”章}。FID只比较两个分布，不检查任何一幅图是否满足它自己的条件。它还依赖Inception特征和样本数：原论文用5万幅生成图像估计均值和协方差，样本少时估计误差大，所以只有在相同样本数、相同特征提取实现下得到的FID才能互相比较。

第二组衡量结果与文字条件的一致程度。\term{CLIP分数}（CLIPScore）用CLIP ViT-B/32分别编码文字$c$和图像$\bm{X}$，取两个向量的余弦相似度\scite{vision-clipscore2021}：
\begin{equation}
 \mathrm{CLIPScore}(c,\bm{X})
 =w\max\bigl(\cos(\bm{e}_{\mathrm{T}}(c),\bm{e}_{\mathrm{I}}(\bm{X})),0\bigr)\text{，}
 \label{eq:vision-clipscore}
\end{equation}
其中$\bm{e}_{\mathrm{T}},\bm{e}_{\mathrm{I}}$是CLIP的文字与图像编码器，原论文取$w=2.5$，只用于拉伸取值范围，并在文字前加上“A photo depicts”。余弦为0.30时分数为0.75，余弦为负时记为0。原论文提出它是为了评价图像描述，生成任务借用它时比较的是同一对象：图文是否相关。它是一个整图对整句的相似度，对数量、左右关系和颜色归属并不敏感，\ref{sec:vision-generation-reranking}节会给出按它排序选错候选的例子。主体保真也用相同的余弦形式：DreamBooth把生成图像与真实主体照片的DINO特征余弦相似度在所有配对上取平均，称为DINO分数，用CLIP图像特征时称为CLIP-I，用图文特征时称为CLIP-T\scite{vision-dreambooth2023}。

第三组依靠人的判断。人工偏好评价把两个系统在同一条件下的结果并排给评审，请他们选出更好的一幅或判为相当。若$n$次比较中系统A胜$n_{\mathrm{w}}$次、平$n_{\mathrm{t}}$次，则A的胜率记为$(n_{\mathrm{w}}+n_{\mathrm{t}}/2)/n$。例如100次比较中胜55次、平10次，胜率为0.6。评审被问到的问题决定了这个数字的含义：“哪幅更真实”和“哪幅更符合文字”是两种不同的偏好，GLIDE就分别就照片真实感和与文字的相符程度作了两两比较\scite{vision-glide2022}。报告人工偏好时，应写明问题措辞、比较次数、评审来源和候选的选取方式。

第四组用于有输入或有真实答案的任务。编辑任务需要分别衡量“该改的是否改了”和“该留的是否留了”。对于明确的保留掩码$\bm{M}\in\{0,1\}^{h\times w}$，$m_{i,j}=1$表示保留区，可以计算保留区的平均变化
\begin{equation}
 E_{\mathrm{keep}}=
 \frac{\sum_{i,j,k}m_{i,j}
       |x'_{i,j,k}-x_{i,j,k}|}
 {c\sum_{i,j}m_{i,j}}\text{，}
 \label{eq:vision-edit-preservation}
\end{equation}
其中$\bm{X}$与$\bm{X}'$是原图与结果，$c$是通道数，分母要求至少有一个保留像素。原样返回原图时$E_{\mathrm{keep}}=0$，但请求根本没有执行，所以这一项必须与修改完成度一起报告。新视角合成可以在留出的测试相机下与真实照片$\bm{Y}$比较。峰值信噪比$\mathrm{PSNR}=10\log_{10}(V_{\max}^2/\mathrm{MSE})$以均方误差为基础，$V_{\max}$是像素取值上限；像素在$[0,1]$内、均方误差为0.001时，PSNR为30\,dB。结构相似性（SSIM）比较局部窗口内的均值、方差和协方差，LPIPS比较深度特征之间的距离，二者的定义与适用条件见\ref{sec:vision-restoration-psnr-ssim}节和\ref{sec:vision-restoration-lpips}节。

这些指标各自只回答一个问题。下面四个小节讨论的困难，正是让这些问题难以同时得到肯定回答的原因。

\subsection{一对多映射带来的条件遵循困难}
\label{sec:vision-generation-challenge-one-to-many}

“雪地里的红色玩具狗”对应无数幅合格的图像：狗可以坐着或站着，天空可以晴或阴。条件与图像之间不是函数关系，一个条件$c$对应一个分布$p(\bm{X}\mid c)$。若像回归问题那样训练一个确定性网络$\widehat{\bm{X}}=G(c)$并最小化平方误差，最优输出是条件均值$\mathbb{E}[\bm{X}\mid c]$。设某个像素在合格图像中一半取0（阴影里）、一半取1（雪面上），平方误差$\frac12(y-0)^2+\frac12(y-1)^2$对$y$求导得$2y-1$，在$y=0.5$处最小，网络会输出一个两种情形都不是的灰色。许多像素同时取平均，就得到模糊、缺少细节的图像。所以生成模型必须表示整个分布并从中采样；而一旦采样，每幅输出都带有随机性，有的满足条件、有的不满足，条件遵循就成了一个概率意义上的要求。

经典方法不去学习一个分布，而是直接从样例中取值。Efros与Leung的纹理合成在样例图中找出所有与当前邻域相似的位置，把它们中心像素的经验分布当作条件分布，再从中随机抽取一个值\scite{vision-efros-leung1999}。因为每次取的都是样例中真实出现过的值，输出保留了样例的锐利纹理，没有平均化。纹理本身可以用局部统计刻画，经典描述见\ref{sec:vision-features-texture}节。不过这种非参数采样只能复制样例中已有的统计规律，读不懂“红色”“雪地”这样的语义条件，它在编辑中的用法见\ref{sec:vision-generation-exemplar-inpainting}节。

学习方法用生成模型表示$p(\bm{X}\mid c)$：生成对抗网络让判别器惩罚不像真实样本的输出，扩散模型学习逐步去噪的条件分布。平均化由此得到解决，问题转成两件事：怎样让采样结果更集中地落在满足条件的区域，以及一次采样不满足时怎么办。\ref{sec:vision-generation-creation}节从条件接口、采样中的引导与采样后的重排、生成后的修订三方面回答；\ref{sec:vision-generation-pix2pix}节的pix2pix则在有配对数据的转换中，用局部判别器抵消$\ell_1$像素误差带来的平均化。

\subsection{组合语义带来的属性绑定错误}
\label{sec:vision-generation-challenge-binding}

“红色玩具狗在蓝色杯子左边”由两个对象、两种颜色和一条空间关系组成。常见的错误是画出蓝色的狗和红色的杯子：每个词都在画面中出现了，但颜色附着到了错误的对象上。\term{属性绑定}（attribute binding）指把每个属性正确地附着到它所修饰的对象上，这里的困难来自组合：对象和属性各自都不难画，难的是按句子规定的方式把它们配对，数量和空间关系也属于同一类组合要求。

以交叉注意力读取文字的生成网络中，每个空间位置都可以按学到的权重读取所有词元；结构上并没有一个机制规定“读取了‘狗’的那些位置也必须读取‘红色’而不是‘蓝色’”，这种配对只能从训练数据中学到。训练数据中少见的组合，例如蓝色的狗，学得往往更不充分。

经典图像合成用图层显式地完成绑定。设计者从一张照片里抠出红色玩具狗，从另一张照片里抠出蓝色杯子，分别放到画布上指定的位置；每个图层的颜色来自它自己的源照片，不会串换，位置也由人直接摆放。这种做法的问题在于拼接处：两张照片的光照和色调不同，直接粘贴会留下明显的边界，\ref{sec:vision-generation-poisson}节的泊松图像编辑就是为消除这种边界而设计的。它的代价是需要现成的源照片和人工摆放，生成不了任何照片里都没有的组合。

学习方法目前没有一个能保证绑定正确的结构，能做的是把绑定要求外显出来，并在生成后逐项检查。\ref{sec:vision-generation-conditions}节用空间图规定布局，用参考角色记录规定每幅参考负责哪个对象；\ref{sec:vision-generation-selection}节先按要求记录筛选候选，再排序；\ref{sec:vision-generation-unireason-refinement}节针对具体错误作一轮修订；\ref{sec:vision-generation-condition-check}节把绑定正确率作为单独的分项报告。

\subsection{局部修改带来的内容保持冲突}
\label{sec:vision-generation-challenge-preservation}

把狗旁边的杯子换成花瓶时，花瓶要立在桌面上，受光方向与周围一致，它的影子会落到杯子原来轮廓之外；杯子原来的影子也应随之消失。若严格要求杯子轮廓以外的像素一个都不变，新旧内容在边界处就会接不上，旧影子也留在原处；若允许外面变化，又可能顺带改动了不该改的狗。保持与修改之间的这种牵制，在转换中以另一种形式出现：照片转成水彩画时每个像素的颜色都要改变，需要保持的是姿态、数量和布局这些结构，而不是像素值。

经典方法的回答是只从已有图像中搬运内容。泊松图像编辑保持区域外的目标图像完全不变，区域内只保留源图像的梯度，并要求结果在区域边界上与周围的值相等，从而消除拼接缝（\ref{sec:vision-generation-poisson}节）。移除杯子之后的空洞，可以用同一幅图中别处的桌面图像块逐块填补，Criminisi等人的方法决定先填哪里（\ref{sec:vision-generation-exemplar-inpainting}节），PatchMatch则让“找到最相似的图像块”这一步快到可以交互使用（\ref{sec:vision-generation-patchmatch}节）。这些方法的局限也来自同一点：新内容只能来自现有图像，粘贴的物体在泊松编辑后整体颜色会被周围带偏，透视和光照方向也不会自动调整，桌面上本来没有的东西更是补不出来。

学习方法能按文字生成原图中不存在的内容。带掩码的修补模型学会根据可见区域生成与之协调的新内容，指令编辑模型不需要掩码，由模型自己判断改哪里，但未提及的区域也可能变化；多轮编辑还会让主体逐轮漂移。转换中，循环一致性和感知损失分别用往返重建和固定网络的特征来保持结构。\ref{sec:vision-generation-translation}节和\ref{sec:vision-generation-editing}节回答这一挑战，\ref{sec:vision-generation-change-check}节讨论怎样分别核验修改与保持。

\subsection{有限观测带来的三维一致性困难}
\label{sec:vision-generation-challenge-3d}

绕玩具狗在两个高度上各拍20张照片，要求给出从正后方偏上看的图像。新图像中的每个像素对应空间中的一条射线，它的颜色必须与所有拍到同一表面点的照片相符；相机绕狗转动时，耳朵应留在同一个三维位置，而不是每换一个视角就重新画一次。二维生成模型可以画出一幅看起来合理的背面，却没有任何机制保证换一个角度后它与前一幅、与真实照片一致。狗的底面从未被拍到，这一部分无论用什么方法都没有观测依据。

经典路线以多视图几何为基础：运动恢复结构从照片中同时估计各相机的位姿和一组稀疏的三维点，稠密重建再得到表面（\ref{sec:vision-formation-multi-view}节）；也可以不显式重建几何，而把照片看成对光线的采样，新视角直接由已采样的光线插值得到，这就是基于图像的渲染。前一条路线在毛绒、半透明和细小结构上难以得到准确的表面，后一条路线需要非常密集的拍摄。

学习方法为每个场景拟合一个可以从任意相机渲染的三维表示，并通过可微渲染让它在所有训练视图下与照片一致。一致性来自所有视图共享同一个三维表示，而不是逐幅图像各自生成。NeRF和3D Gaussian Splatting是两种代表，见\ref{sec:vision-generation-novel-view}节；没有拍到的表面仍然没有观测支持，\ref{sec:vision-generation-costs}节讨论怎样在交付时标明这一点，需要靠生成先验补全的情形留给\ref{chap:vision-multimodal-generation}章。

四个挑战并不一一对应四类任务，而是按约束由弱到强逐步叠加。创建只受条件约束，一对多和属性绑定是它的主要困难；转换和编辑在此之上多了一幅输入，保持与修改的冲突随之出现，转换在整幅图上保持结构，编辑只在局部修改；新视角合成要求结果由真实观测支持，三维一致性取代了“看起来合理”成为首要标准。所以本章按创建、转换、编辑、新视角合成的次序展开，每一类任务都继承前一类的困难；最后一节再按条件遵循、输入保持、观测支持与成本三个维度汇总核验方法，它们分别对应前两个挑战、第三个挑战和第四个挑战。

\section{新图像的创建}
\label{sec:vision-generation-creation}

创建新图像时，输入是一组条件：一段文字，可能还有姿态图、深度图或几张主体照片；输出是一幅或几幅满足条件的图像。难点有三处。第一，文字很难精确规定位置和形状，几张照片又难以通过文字表达，因此条件要按其约束的对象选用不同接口送入模型。第二，扩散模型的采样是随机的，同一组条件会给出质量和条件满足程度不同的候选，需要决定在采样中修正还是在采样后挑选。第三，生成后才能看到的错误，例如数量多了一只狗，需要一个“检查—修订”的环节。这三处难点对应创建流程中依次进行的步骤：先把请求整理成可检查的要求，再组织条件，然后生成并选择候选，最后在必要时作一次有记录的修订。这些步骤是任务流程中的职责划分，不要求模型内部恰好有四个同名模块。

\subsection{创作要求的确定}
\label{sec:vision-generation-requirements}

生成模型接受的是一段提示，验收者却需要一份可以逐项打勾的清单。考虑请求“红色玩具狗放在蓝色杯子左边”。它包含两个对象、各自的颜色、每种对象的数量以及一条左右关系。把它写成记录，就是“玩具狗，一只，红色；杯子，一只，蓝色；狗在杯子左边”。这份记录供后续生成选择和结果核验使用，并不假定生成器内部也存有同样的对象表。

把请求拆开以后，才能看清几层不同的要求。对象出现和属性绑定是两层：画面里同时有红色和蓝色，不说明红色属于狗、蓝色属于杯子；画面里有狗也有杯子，不说明数量和左右关系正确。若还要求杯子上写“OPEN”，字符应作为单独一项检查，因为整体图文相似度对几个字母的拼写并不敏感。

输出用途还决定哪些部分可以变化。插图允许材质、光照和背景有多种合理选择；产品展示要求主体形状固定；排版素材还要规定留白和裁剪范围。条件没有指定的部分可以自由变化，已经指定的部分不能因为“画面自然”而放宽。一个请求通常有许多可接受结果，因此要求记录应写出必须满足的项和允许变化的项，而不是把某一幅参考图的全部像素当作唯一答案。

\subsection{创作条件的组织}
\label{sec:vision-generation-conditions}

有了要求记录，下一步是把它变成模型能读取的条件。玩具狗例子里的要求分属三类信息：“雪地、杯子、红色”这类语义适合用文字表达；“狗抬起左前爪、坐在画面左侧”这类几何关系，用文字描述既冗长又不精确；“就是照片里这只玩具狗”这类身份信息，文字几乎无法表达，只能从参考照片中取得。条件组织要解决的问题是：每类信息用什么形式编码，在生成过程的哪一步进入，以及几类条件冲突时怎样判断结果。

按条件的形式，本节分四种情形。文字条件下的生成直接使用\ref{chap:vision-multimodal-models}章的LDM和unCLIP，条件经文字编码器进入去噪网络，不需要额外训练；空间图条件下的生成需要把一幅与输出对齐的图接入预训练模型，ControlNet用一条可训练的复制支路完成这件事，原模型参数不动；单一主体参考下的生成需要让模型记住一个具体主体，DreamBooth把主体信息写进模型参数本身；多参考图像下的生成则要再回答每幅参考负责什么。从文字到多参考，条件携带的信息越来越具体，代价也从“只改提示”增加到“训练支路”“微调整个模型”和“需要显式规定参考角色”。

\subsubsection{文字条件下的生成}
\label{sec:vision-generation-latent-entry}

只有文字时，最直接的工具是\ref{sec:vision-models-ldm}节的LDM。它先用自编码器把图像压缩成空间潜网格，在潜空间中训练以文字为条件的去噪网络，文字经编码器形成词元序列，再通过交叉注意力被每个空间位置读取；采样时从高斯噪声出发反复去噪，最后由解码器还原像素\scite{pgm-generative-rombach2022}。\ref{sec:vision-models-unclip}节的unCLIP则先由先验模型把文字映射为CLIP图像向量，再由解码器根据这个向量生成图像；给定一幅参考图时，也可以直接编码它的CLIP向量，跳过先验，生成语义相近的变体\scite{vision-unclip2022}。两者的模型结构和训练在该章已讲，这里只讨论使用时的输入组织、超参数和失败情形。

使用LDM时，一次生成由以下量共同决定：提示文本、目标分辨率、随机种子、采样器与步数、引导强度。分辨率先换算成潜网格大小。若自编码器在高宽方向各压缩8倍，$512\times512$的目标对应$64\times64$的潜网格，通道数由所选自编码器规定，例如4。种子决定初始高斯噪声$\bm{z}^{(K)}\in\mathbb{R}^{64\times64\times4}$，采样器按噪声日程调用去噪网络$K$次，每次读取文字特征，把当前状态更新为下一状态；最后解码得到$\bm{X}\in\mathbb{R}^{512\times512\times3}$，再按训练时的像素归一化反变换到显示范围。潜数组不是可交付的RGB图片；改变长宽比也必须相应改变潜网格，只在最后拉伸图像会让采样条件与报告不符。

unCLIP的参考图路径与LDM的空间潜网格不同。CLIP图像向量是一个整图汇总向量，参考图经它进入解码器时，被保留的是“这是一只红色毛绒玩具狗”一类语义，耳朵的精确形状、身上的标志和原背景并不保证保留。所以unCLIP适合产生与参考风格或语义相近的多个变体，不适合要求“同一只玩具”的身份保持。使用时需要同时保存先验、解码器、上采样器及其预处理版本，只保留CLIP编码器并不能重现生成路径。

文字条件的常见失败有三类。第一，属性绑定错误，例如狗被画成蓝色、杯子被画成红色，原因是交叉注意力并没有为“红色”指定一个对象框，词语与区域的对应完全靠训练数据学到。第二，数量与空间关系不准，例如“两个杯子”画成三个。第三，小字和细小结构在潜空间压缩中丢失。这些失败不能靠更长的提示保证消除，后文的空间条件、候选选择和修订分别针对其中一部分。实际比较时应固定模型、提示、分辨率和候选数，例如对同一描述用四个种子生成，观察姿态、背景和细节的变化；相同种子只有在模型、采样器和实现都一致时才可复现。

\subsubsection{空间图条件下的生成}
\label{sec:vision-generation-controlnet}

文字“狗抬起左前爪，坐在画面左三分之一处”很难让LDM把爪子放到准确位置。若手头有一幅姿态关键点图、边缘图或深度图，它与输出图像逐像素对齐，直接规定了布局。问题是，已有的文字生成模型只有文字条件接口，没有读取空间图的入口。从头训练一个同时读取文字和空间图的模型代价很高；直接在已有模型上微调全部参数，又可能在只有几万对条件图的数据上破坏原模型学到的生成能力。ControlNet解决的就是“在不损坏预训练模型的前提下，为它增加空间条件接口”\scite{vision-controlnet2023}。

ControlNet的主要想法是：冻结原模型，把原模型编码部分复制一份作为可训练支路，由支路读取空间条件，再通过初始化为零的卷积把支路输出加回原模型。复制保证支路从一开始就具备原模型的特征提取能力，而不是从随机权重学起；零初始化保证训练开始时支路的输出为零，原模型的行为与微调前完全相同，训练只会逐渐加入控制信号。\term{零卷积}（zero convolution）指权重和偏置都初始化为零的$1\times1$卷积，它是支路与主网络之间的唯一连接。

先看一个模块上的计算。设预训练模块为$F_{\bm{\theta}}$，复制模块为$F_{\bm{\theta}_{\mathrm{c}}}$，输入特征为$\bm{h}$，条件特征为$\bm{c}$，两个零卷积为$Z_1,Z_2$，则模块输出为
\begin{equation}
 \bm{h}'=F_{\bm{\theta}}(\bm{h})+
 Z_2\!\left(F_{\bm{\theta}_{\mathrm{c}}}
       \bigl(\bm{h}+Z_1(\bm{c})\bigr)\right)\text{。}
 \label{eq:vision-controlnet}
\end{equation}
$Z_1$把条件特征注入复制模块的输入，$Z_2$把复制模块的输出加回主路径。训练开始时两者输出都为零，$\bm{h}'=F_{\bm{\theta}}(\bm{h})$。零输出并不意味着支路没有梯度：$Z_2$的权重梯度等于上游误差与它的输入$F_{\bm{\theta}_{\mathrm{c}}}(\cdot)$的乘积，后者不为零，所以$Z_2$第一次更新后就变成非零，梯度随后才能继续传到复制模块和$Z_1$。图\ref{fig:vision-controlnet}标出了两条路径的关系。

\begin{figure}[htbp]
 \centering
 % Single-module sketch of ControlNet Eq.(2); no complete U-Net is implied.
\begin{tikzpicture}[x=1mm,y=1mm,font=\figurefont\small,
 b/.style={draw=BookBlue,fill=BookBlue!10,line width=.8pt,
  minimum width=16mm,minimum height=9mm,inner sep=1mm,align=center},
 op/.style={circle,draw=BookInk,fill=white,line width=.8pt,
  minimum size=5mm,inner sep=0},
 a/.style={-{Stealth[length=2mm]},BookInk,line width=.9pt,
  rounded corners=1mm}]
 \node at (5,0) (h) {$\bm{h}$};
 \node[b,draw=BookMuted,fill=BookPaper,minimum width=25mm] at (42,0)
  (base) {$F_{\bm{\theta}}$};
 \node[op] at (82,0) (sum) {$+$};
 \node at (101,0) (out) {$\bm{h}'$};
 \node[op] at (30,-25) (cond) {$+$};
 \node[b,minimum width=23mm] at (55,-25) (copy) {$F_{\bm{\theta}_{\mathrm{c}}}$};
 \node[b,minimum width=12mm] at (82,-25) (z2) {$Z_2$};
 \node at (5,-45) (c) {$\bm{c}$};
 \node[b,minimum width=12mm] at (30,-45) (z1) {$Z_1$};
 \draw[a] (h)--(base);
 \draw[a] (base)--(sum);
 \draw[a] (sum)--(out);
 \draw[BookInk,line width=.9pt] (14,0)--(14,-25);
 \draw[a] (14,-25)--(cond);
 \fill[BookInk] (14,0) circle (.6mm);
 \draw[a] (cond)--(copy);
 \draw[a] (copy)--(z2);
 \draw[a] (z2)--(sum);
 \draw[a] (c)--(z1);
 \draw[a] (z1)--(cond);
 \node[anchor=south,text=BookMuted] at (42,7) {冻结};
 \node[anchor=north,text=BookBlue] at (55,-32) {可训练副本};
 \node[anchor=north,text=BookBlue] at (82,-32) {从零连接};
\end{tikzpicture}

 \caption{ControlNet的模块级教学示意。主模块参数$\bm{\theta}$冻结，复制模块及零初始化连接可训练；相同输入进入两侧，条件只通过支路注入。零初始化让控制输出从零开始，训练后经加法调节主路径。完整网络在多个尺度接入控制特征，图中仅展开一处。}
 \label{fig:vision-controlnet}
\end{figure}

零卷积启动的顺序可以手算。把外层连接简化为标量$a$乘以支路输出$q$，主路径输出为$f$，目标为$y$，损失为$(f+aq-y)^2/2$。取$a=0,q=2,f=0.2,y=0.6$，则$\partial L/\partial a=(f+aq-y)q=-0.8$，而$\partial L/\partial q=(f+aq-y)a=0$。以学习率0.1更新一次得到$a=0.08$；下一步$\partial L/\partial q$就不再为零。这个教学例只说明梯度怎样逐层启动，不表示完整网络的收敛步数。

把这个模块放回整个网络。原论文以Stable Diffusion为基础，它的去噪U-Net在$64\times64\times4$的潜网格上工作，由编码部分、中间块和带跳连的解码部分组成；编码与解码部分各有12个块，分布在$64,32,16,8$四种分辨率上，每种分辨率3个块。ControlNet复制这12个编码块和1个中间块，复制支路的13个输出各经一个零卷积，分别加到U-Net解码部分的12条跳连和中间块的输出上；解码部分本身不复制，保持冻结。论文没有列出各块的通道数，表\ref{tab:vision-controlnet-structure}的通道与层内组成取自Stable Diffusion v1公开的配置文件：基础通道320，四个分辨率的通道倍率为$1,2,4,4$，前三个分辨率的块含空间Transformer，文字由冻结的CLIP文字编码器编码为77个768维词元。

\begin{table*}[tbp]
 \centering
 \small
 \begin{tabularx}{\linewidth}{@{}p{30mm}Xp{24mm}p{22mm}p{12mm}p{20mm}@{}}
  \toprule
  阶段或层组 & 层类型 & 卷积核与步长 & 输出通道 & 重复次数 & 输出尺寸\\
  \midrule
  条件编码器（论文） & 卷积，ReLU & $4\times4$，步长2 & 16、32、64、128 & 4 & 论文称$64\times64$\\
  条件编码器（官方代码） & 卷积，SiLU；末层为零卷积 & $3\times3$，其中3层步长2 & 16、16、32、32、96、96、256，末层320 & 8 & $64\times64$\\
  输入卷积 & 卷积 & $3\times3$，步长1 & 320 & 1 & $64\times64$\\
  编码分辨率1 & 残差块，空间Transformer & $3\times3$，步长1 & 320 & 2 & $64\times64$\\
  下采样 & 卷积 & $3\times3$，步长2 & 320 & 1 & $32\times32$\\
  编码分辨率2 & 残差块，空间Transformer & $3\times3$，步长1 & 640 & 2 & $32\times32$\\
  下采样 & 卷积 & $3\times3$，步长2 & 640 & 1 & $16\times16$\\
  编码分辨率3 & 残差块，空间Transformer & $3\times3$，步长1 & 1280 & 2 & $16\times16$\\
  下采样 & 卷积 & $3\times3$，步长2 & 1280 & 1 & $8\times8$\\
  编码分辨率4 & 残差块 & $3\times3$，步长1 & 1280 & 2 & $8\times8$\\
  中间块 & 残差块，空间Transformer，残差块 & $3\times3$，步长1 & 1280 & 1 & $8\times8$\\
  零卷积（仅ControlNet） & 卷积 & $1\times1$，步长1 & 与所接的块相同 & 13 & 与所接的块相同\\
  解码部分（冻结，不复制） & 拼接跳连后的残差块（除$8\times8$外含空间Transformer），分辨率之间上采样 & $3\times3$，步长1 & 1280、1280、640、320 & 每个分辨率3个块 & $8\to64$\\
  输出层 & 组归一化，SiLU，卷积 & $3\times3$，步长1 & 4 & 1 & $64\times64$\\
  \bottomrule
 \end{tabularx}
 \caption{ControlNet所依附的Stable Diffusion v1去噪U-Net与新增部分。编码部分从“输入卷积”到“编码分辨率4”共12个块，ControlNet把它们和中间块各复制一份；“零卷积”一行是新增的13个连接。条件编码器给出论文文字与官方代码两种描述，二者不一致，见正文。通道与层内组成取自公开配置，论文正文只给出块数与分辨率。}
 \label{tab:vision-controlnet-structure}
\end{table*}

表中的两种模块按以下顺序计算。残差块对输入依次作组归一化、SiLU、$3\times3$卷积，加上由噪声步编码而来的向量，再作组归一化、SiLU、$3\times3$卷积，最后与输入相加（通道数变化时输入先经$1\times1$卷积）。空间Transformer块把$h\times w$个位置展平为词元序列，依次经过自注意力、读取文字词元的交叉注意力和前馈层，每一步都带残差连接，再恢复为特征图。文字条件只经交叉注意力进入，冻结主网络和复制支路都读取同一份文字词元。

空间条件图的分辨率是$512\times512$，而支路在$64\times64$的潜网格上工作，因此先由一个小卷积网络$E$把条件图编码成$64\times64$的特征，再经$Z_1$加到复制支路输入卷积的输出上。论文文字说$E$由4层$4\times4$、步长2的卷积组成，通道依次为16、32、64、128，激活为ReLU；但4次步长为2的下采样会把512降到32，与它所说的64不符。作者公开的实现用8层$3\times3$卷积，其中3层步长为2，激活为SiLU，最后一层是输出320通道的零卷积，恰好把512降到64并与输入卷积的通道数对齐。复现时应以所用代码为准，并在报告中写明。数据流可以概括为：带噪潜变量与文字进入冻结U-Net；同一带噪潜变量经复制的输入卷积后加上编码后的条件图，进入复制编码器；复制编码器各块的输出经零卷积加到冻结U-Net解码端的相应位置；最后由冻结U-Net输出噪声预测。论文报告，在单张40GB A100上，带ControlNet的训练每次迭代比单独微调Stable Diffusion多用约23\%的显存和34\%的时间，因为冻结部分不需要计算参数梯度。

训练数据是“目标图像、文字、条件图”三元组，条件图通常由检测器从目标图像自动提取，例如用边缘检测得到边缘图、用姿态估计得到关键点图。训练目标仍是扩散模型的噪声预测误差。设目标图像的潜表示为$\bm{z}_i$，条件图为$\bm{C}_i$，提示为$c_i$；在小批量中为每个样本抽噪声步$k_i$和噪声$\bm{\epsilon}_i$，构造带噪状态$\bm{z}_i^{(k_i)}$，则
\begin{equation}
 L_{\mathrm{ctrl}}=
 \frac{1}{m}\sum_{i=1}^{m}
 \left\lVert\bm{\epsilon}_i-
 \epsilon_{\bm{\theta},\bm{\psi}}
  (\bm{z}_i^{(k_i)},k_i,c_i,\bm{C}_i)
 \right\rVert_2^2\text{，}
 \label{eq:vision-controlnet-training}
\end{equation}
其中$m$是小批量大小，$\bm{\psi}$汇集复制支路、条件编码器和零卷积的参数，只有$\bm{\psi}$被更新，$\bm{\theta}$保持冻结。原论文训练时以50\%的概率把文字替换为空串，迫使模型从条件图本身读出语义，而不是只依赖文字。论文还报告了一种“突然收敛”现象：模型在一段训练中几乎不遵循条件，随后在不到一万步内突然开始遵循。每种条件类型（边缘、深度、姿态等）分别训练一个支路。

推理时同时给出文字和条件图，在每个去噪步注入支路特征。支路输出可以乘以一个控制强度系数：强度为零时退回原模型，强度增大时输出更贴近条件图的结构。多种条件可以同时使用，原论文把多个ControlNet的输出直接相加。但若深度图规定前爪在身体后方、姿态图又规定前爪伸向相机，两种条件没有一致的几何解，相加并不能消解矛盾。常见失败还包括条件提取器本身出错，例如姿态估计把左右爪弄反，ControlNet会忠实地画出错误姿态。因此需要先核验条件图，再记录支路版本、条件分辨率和控制强度。空间条件只规定布局，不规定“这是哪一只玩具狗”，后者需要下面的主体适配。

\subsubsection{单一主体参考下的生成}
\label{sec:vision-generation-dreambooth}

文字“玩具狗”指一个类别，ControlNet只规定布局，二者都无法让模型画出照片里那只特定的玩具狗，例如它左耳的缝线和背上的补丁。unCLIP的图像向量又压缩掉了这些细节。要让模型在新场景、新姿态中重现一个具体主体，就需要把这个主体的外观写进模型。DreamBooth的做法是用该主体的3到5张照片微调整个文字生成模型，并把主体绑定到提示里的一个特殊标识上\scite{vision-dreambooth2023}。

主要想法分两部分。其一，选择一个在词表中很少出现、本身几乎没有语义的词元作为标识，记作“[V]”，训练提示写成“a [V] toy dog”。因为标识原本没有含义，微调可以把主体外观绑定到它上面，同时类别词“toy dog”提供了类别先验，让模型不必从零学习“玩具狗长什么样”。其二，只在几张照片上微调会带来两种退化：模型会把“toy dog”这个类别词也缩窄成这一只玩具，即\term{语言漂移}（language drift）；输出还会趋同于几张训练照片的姿态和视角。\term{类别先验保持}（class-specific prior preservation）用来抑制这两种退化：在微调前，先用原模型以提示“a toy dog”生成约1000幅一般玩具狗图像，微调时让模型继续能生成这些图像，从而保持类别词的原有含义。

DreamBooth不改变模型结构，输入输出都与原文字生成模型相同：输入是提示与带噪图像（或潜表示），输出是去噪预测。所以它的“结构”就是所微调的基础模型的结构，需要交代的是微调了其中哪些部分。原论文在Imagen和Stable Diffusion上实验，二者的组成不同。Imagen是级联模型：一个以T5-XXL文字编码为条件、输出$64\times64$图像的基础扩散模型，接两个文字条件的超分辨率扩散模型，分别完成$64\to256$和$256\to1024$。原论文发现，主体保真的最好结果来自微调全部层，包括以文字编码为条件的层；只微调基础模型时，未经微调的超分辨率模型会给主体添上错误的高频纹理，所以$64\to256$的超分辨率模型对多数主体都必须一起微调，$256\to1024$的模型只对细节特别多的主体有帮助。超分辨率模型训练时会给低分辨率输入加噪声并以噪声强度为条件，称为噪声条件增强；原论文在微调$64\to256$模型时把这一噪声强度从训练时的$10^{-3}$降到$10^{-5}$，否则主体的细密花纹会被抹平。Stable Diffusion只有一个在潜空间工作的去噪U-Net（结构见表\ref{tab:vision-controlnet-structure}），微调的是这个U-Net，文字编码器可以一并微调，自编码器的解码器保持不变。标识的选择需要经过分词器检查：原论文对Imagen从T5-XXL词表中编号5000到10000、长度不超过3个字符的词元中随机选取，随意拼出的一串字母可能被切成多个常用词元，反而带有已有含义。

训练目标由两项组成：
\begin{equation}
 L=L_{\mathrm{subj}}+\lambda_{\mathrm{pr}}L_{\mathrm{class}}\text{，}
 \label{eq:vision-dreambooth}
\end{equation}
原论文取$\lambda_{\mathrm{pr}}=1$。原论文以预测干净图像的形式写去噪目标。设主体照片为$\bm{x}$，其提示为$c_{\mathrm{subj}}$（“a [V] toy dog”），加噪状态为$\alpha_k\bm{x}+\sigma_k\bm{\epsilon}$，主体项为
\begin{equation}
 L_{\mathrm{subj}}=
 \mathbb{E}_{\bm{x},k,\bm{\epsilon}}
 \left[
 w_k\left\lVert
 f_{\bm{\theta}}(\alpha_k\bm{x}+\sigma_k\bm{\epsilon},
                  k,c_{\mathrm{subj}})-\bm{x}
 \right\rVert_2^2\right]\text{，}
 \label{eq:vision-dreambooth-subject}
\end{equation}
其中$\alpha_k,\sigma_k$是噪声日程系数，$w_k$是噪声步权重，$f_{\bm{\theta}}$输出对干净图像的估计。类别项把$(\bm{x},c_{\mathrm{subj}})$替换为预先生成的类别图像$\bm{x}_{\mathrm{pr}}$与普通提示$c_{\mathrm{pr}}$（“a toy dog”），另抽噪声和噪声步。主体项教模型“[V]是这一只”，类别项教模型“toy dog仍然是一般玩具狗”，二者监督不同的提示。若基础模型输出噪声预测，可由$\hat{\bm{x}}=(\bm{z}^{(k)}-\sigma_k\hat{\bm{\epsilon}})/\alpha_k$换算，得$\lVert\hat{\bm{x}}-\bm{x}\rVert_2^2=(\sigma_k^2/\alpha_k^2)\lVert\hat{\bm{\epsilon}}-\bm{\epsilon}\rVert_2^2$，所以改写成噪声误差时要同步变换噪声步权重。

训练时每个小批量同时抽主体照片和类别图像，分别计算两项后加权求和，反传更新全部被微调的参数。两种基础模型上的训练配置见表\ref{tab:vision-dreambooth-config}。论文没有给出批大小和优化器设置。训练结束后保存微调权重、标识和类别提示。推理时输入“a [V] toy dog sitting in the snow”，即可在雪地中生成这只玩具狗，不再需要参考照片。

\begin{table}[htbp]
 \centering
 \small
 \begin{tabularx}{\linewidth}{@{}p{22mm}XX@{}}
  \toprule
  项目 & Imagen & Stable Diffusion\\
  \midrule
  基础模型组成 & $64\times64$基础模型，$64\to256$与$256\to1024$两级超分辨率 & 潜空间去噪U-Net，自编码器，文字编码器\\
  微调的部分 & 全部层，含读取文字编码的层；$64\to256$超分辨率模型一并微调，$256\to1024$可选 & 去噪U-Net，文字编码器可选；解码器不变\\
  主体照片 & 3到5张 & 3到5张\\
  类别图像 & 约1000幅，也可更少 & 约1000幅，也可更少\\
  $\lambda_{\mathrm{pr}}$ & 1 & 1\\
  学习率 & $10^{-5}$ & $5\times10^{-6}$\\
  迭代次数 & 约1000 & 约1000\\
  训练时间 & 单个TPUv4约5分钟 & 单张A100约5分钟\\
  其他 & 标识取自T5-XXL词表编号5000到10000、不超过3个字符的词元；超分辨率噪声条件增强由$10^{-3}$降至$10^{-5}$ & 论文未给出标识的选取范围\\
  \bottomrule
 \end{tabularx}
 \caption{DreamBooth在两种基础模型上的微调配置，数值取自原论文。DreamBooth不改动网络结构，表中列出的是被微调的部分与训练设置；批大小与优化器论文未给出。}
 \label{tab:vision-dreambooth-config}
\end{table}

DreamBooth的边界也来自它的训练数据。若只有正面照片，生成背面时的纹理来自类别先验，而不是这只玩具；增大$\lambda_{\mathrm{pr}}$保持了类别多样性，却不能补出从未拍到的外观。训练过久会让模型复制参考照片的背景；多个主体各自微调后再组合，又会出现标识绑定混淆。原论文用DINO特征相似度和CLIP图像相似度衡量主体保真，用CLIP图文相似度衡量提示遵循，二者需要同时报告，因为单纯复制参考照片会让前者很高、后者很低。

\subsubsection{多参考图像下的生成}
\label{sec:vision-generation-reference-roles}

现在请求变成：“让第一张照片里的玩具狗，坐在第二张照片里那种蓝釉杯子旁边，画成第三张图的水彩风格”。三幅参考分别负责主体身份、局部外观和整体风格。上面三种情形都只有一种条件，此时新出现的难点是参考之间的职责分配：若只写“参考这三张图”，模型可能把杯子的蓝色涂到狗身上，把第一张的桌面背景一起复制，或者把水彩风格只用在杯子上。

解决这个问题的第一步不在模型里，而在要求记录中：为每幅参考写明它对应的目标对象和应取用的属性，例如“参考1：玩具狗的身份（耳朵、斑纹）；参考2：杯子的釉色与形状；参考3：全图笔触与色调；布局由文字规定”。我们把这种对应称为参考角色绑定。有了它，生成后才能分项检查是否发生串换，例如狗是否变成蓝色、杯子是否保留了形状。

第二步是选用支持多图输入的模型接口，并确认它确实能区分各参考的角色。可选的做法分两类：像DreamBooth那样为主体单独微调，再用其他条件提供其余参考；或使用在推理时直接读取参考图像的接口。\ref{sec:vision-generation-kontext-history}节要讲的FLUX.1 Kontext把参考图像编码为词元序列，与待生成图像的词元一起送入同一个Transformer，并为不同来源的图像分配不同的位置编号；但其原论文主要以单张上下文图像训练，多张参考同时输入不属于它已经验证的能力。所以角色绑定是一份要求规格，不能为只接收单图的模型凭空增添多参考能力。

冲突应在采样前暴露。若参考2是蓝杯子，提示却要求杯子是红色，就要在记录里写明“取杯型、改颜色”；若两处条件要求杯子分别在狗的左右两侧，必须先选定一个。保留属性和可变属性分别记录以后，才能判断输出中的颜色变化是按要求修改还是串换。参考未展示的属性仍由生成先验补充，精确文字仍需读取核验，整体参考相似度不能代替这些分项判断。

四种条件的比较见表\ref{tab:vision-generation-condition-comparison}。固定任务为“在雪地里画这只玩具狗”，它们的差别在于条件提供的信息、模型要付出的适配成本以及典型的失败方式。文字条件不需要任何训练，但只能约束语义；ControlNet用每类条件一次支路训练换来像素级布局控制；DreamBooth用每个主体一次全模型微调换来身份保持；多参考接口省去逐主体训练，却要求模型确实支持多图角色区分，并把角色规格交给使用者写明。

\begin{table}[htbp]
 \centering
 \small
 \begin{tabularx}{\linewidth}{@{}l X X X X@{}}
  \toprule
  条件情形 & 约束的信息 & 需要的适配 & 每次使用的输入 & 典型失败\\
  \midrule
  文字（LDM、unCLIP） & 语义、类别、风格 & 无 & 提示 & 属性串换、数量与位置不准\\
  空间图（ControlNet） & 逐像素对齐的布局 & 每类条件训练一条支路，原模型冻结 & 提示与条件图 & 条件图错误被忠实画出，多条件矛盾\\
  单一主体（DreamBooth） & 具体主体外观 & 每个主体微调整个模型 & 含标识的提示 & 复制背景、未见外观来自类别先验\\
  多参考图像 & 各参考的指定属性 & 依所用接口而定 & 提示、多幅参考与角色记录 & 属性在参考间串换\\
  \bottomrule
 \end{tabularx}
 \caption{四种创作条件在同一创建任务下的比较。“需要的适配”指为得到该能力而额外进行的训练；推理成本另随采样步数和候选数变化。}
 \label{tab:vision-generation-condition-comparison}
\end{table}

\subsection{候选图像的生成与选择}
\label{sec:vision-generation-selection}

条件组织好以后，采样会给出一幅具体图像，但不保证它满足要求记录。固定条件和种子，问题变为怎样提高最终交付图像满足条件的概率。有两种时机可以干预。一是在采样过程中修正每一步的去噪方向，使轨迹更偏向条件，这改变了候选本身；二是先生成多个候选，再用打分器挑选，这不改变任何候选，只改变交付哪一个。两者都可能用到图文相似度，但作用的对象、时机和成本不同，下面依次讨论。

\subsubsection{采样中的引导修正}
\label{sec:vision-generation-glide-guidance}

\ref{sec:vision-models-glide}节已经推导了GLIDE中的两种采样引导\scite{vision-glide2022}。无分类器引导在每一步分别计算同一网络的有条件和无条件噪声预测，按式\eqref{eq:vision-glide-cfg-interface}以强度$g$外推；CLIP引导则用一个在带噪图像上训练过的CLIP模型，按式\eqref{eq:vision-glide-clip-interface}把图文匹配分数对当前图像的梯度乘以强度$g_{\mathrm{clip}}$加到反向均值上。这里不再重复推导，只讨论使用时最常遇到的三个问题：强度怎样选择，代价多大，以及怎样核验它的效果。

先看强度对结果的影响。$g=1$等于不加引导，$g>1$时每一步都沿“有条件减无条件”的方向多走一段。GLIDE报告，提高无分类器引导强度会让样本更贴合文字、画面更像照片，同时样本多样性下降，二者此消彼长；在人工比较中，无分类器引导的结果在真实感和图文相符两方面都比CLIP引导更受偏好。原因之一是CLIP引导沿着一个固定判别模型的梯度修改像素，生成器可能找到让该CLIP模型打高分、人看来却不合理的图像，这类似于针对该模型的对抗样本。强度过大时还会出现颜色过饱和、对比过强和构图单调。所以强度没有“越大越好”的方向，通常在一个中等范围内按任务确定。

强度的选择应作为一次有记录的实验。固定模型、提示、种子集合和采样步数，例如对“红色玩具狗在蓝色杯子左边”用8个种子，分别取$g=1,3,5,7.5,10$生成，得到40幅图。每个强度下按要求记录逐项统计满足率（颜色归属、数量、左右关系），同时统计同一强度下8幅图之间的差异，例如背景和姿态是否趋同。若满足率在某个强度之后不再上升而多样性继续下降，更大的强度就只带来代价。选定的强度只对这组模型、采样器和任务有效，换一种采样器或换成编辑任务都要重新确定。

再看成本。无分类器引导在每一步需要两次网络前向，一次有条件、一次无条件，可以拼成一个批次计算，计算量约为不加引导的两倍。CLIP引导除了去噪网络，每一步还要前向并反传一个噪声感知CLIP模型，而且这个模型需要单独训练，普通CLIP只在清晰图像上训练过，直接接到带噪状态上不能沿用GLIDE的结论。报告引导效果时，应把每步网络调用次数和额外模型一并写明。

核验时要记住，引导提高的是模型自身的条件概率，或者某个CLIP模型的匹配分数，不是要求记录中的每一项。更强的引导可能让“红色”更显著，却同时让狗与杯子的颜色更容易串换。因此引导强度的效果只能由逐项核验判断，不能以平均图文相似度上升代替。

\subsubsection{生成后的候选重排}
\label{sec:vision-generation-reranking}

重排发生在候选已经形成之后。对$k$幅候选$\bm{X}_1,\ldots,\bm{X}_k$计算打分并排序，它不修改任何候选的像素，也不能从全部错误的候选中造出正确答案。最常用的打分是CLIP图文相似度：把每个候选和提示分别编码成单位向量，取内积。GLIDE对照的DALL-E就采用了这种生成后的CLIP重排\scite{vision-glide2022}；不同系统使用的候选总数不同，比较时必须连同选择协议一起阅读。

全局图文相似度只衡量粗粒度的相关性，按它排序会选错。表\ref{tab:vision-generation-candidates}给出四个构造候选：若只取相似度最高者，会选中出现两只狗的A；若先按要求记录逐项筛选，再在通过者中按质量排序，只有D满足全部要求。这里的分项检查由人工或独立的检测、读取工具完成，不是CLIP分数提供的保证。

\begin{table}[htbp]
 \centering
 \small
 \begin{tabularx}{\linewidth}{@{}c c X c@{}}
  \toprule
  候选 & 构造相似度 & 按请求核验的主要结果 & 保留\\
  \midrule
  A & 0.34 & 出现两只玩具狗 & 否\\
  B & 0.32 & 狗为蓝色，杯子为红色 & 否\\
  C & 0.30 & 属性正确，但狗在右侧 & 否\\
  D & 0.28 & 数量、颜色归属与左右关系均满足 & 是\\
  \bottomrule
 \end{tabularx}
 \caption{全局相似度与逐项条件检查的教学对照。分数人为设置，用于展示两种选择规则的区别；不代表任何模型的实测排序。}
 \label{tab:vision-generation-candidates}
\end{table}

由此得到一个两阶段的选择流程：先用硬条件筛选，例如用检测器数狗和杯子的个数、读取它们的位置和颜色；再在通过者中按图文相似度或审美分数排序。重排的成本与候选数成正比：生成$k$幅图就要付出$k$次完整采样，再加上$k$次打分。若一个请求生成了100幅图、人工选出1幅，报告时应写明100幅的生成成本和筛选规则；按“最佳候选”计分与按“一次随机生成”计分，衡量的是两个不同的系统。没有候选通过硬条件时，应返回“条件未满足”，或在明确预算内继续生成或修订，不能把最高分候选自动称为成功。

两种干预的关系可以这样概括。引导在采样中改变每个候选，代价是每步多一次前向（或一次CLIP反传），影响所有候选，并以多样性为代价；重排不改变候选，代价随候选数线性增加，而且只能在已有候选中选择。二者可以组合，例如以中等引导强度生成8个候选，再按硬条件筛选。比较不同系统时应固定引导强度和候选总数，否则无法判断差异来自模型还是选择过程。

\subsection{生成结果的反思修订}
\label{sec:vision-generation-unireason-refinement}

引导和重排都在“一次生成”内起作用。若8个候选都把杯子画在了狗的左侧，重排无能为力，提高引导强度也未必能修正一条具体的空间关系。此时更直接的办法是看着已有的草图，指出哪里不符合要求，再按这条意见修改。生成后的反思修订就是这样一个“生成—检查—编辑”的环节。\ref{sec:vision-models-unireason}节介绍的UniReason 1.0把它训练进一个统一模型：模型先生成草图，再读取草图与原请求，用文字写出偏差说明，最后以草图和偏差说明为条件生成修订图像\scite{vision-unireason-2026}。它的结构、训练数据和联合目标都在那一节讲过，这里只讨论在创作中怎样使用：编辑请求如何形成，修订预算如何规定，修订结果如何验收。

编辑请求的质量决定修订能否成功。偏差说明必须依据实际草图指出“哪个对象、哪个属性或哪条关系需要怎样改变”，并说明哪些部分应保持。例如草图出现两只狗、杯子在狗左侧，合格的请求是“只保留一只红色玩具狗；把蓝杯子移到狗的右侧；保持狗的斑纹和背景”；而“请更符合原提示”只是重复原要求，模型无从知道改哪里。若草图数量和位置都正确，只有阴影方向与光源矛盾，请求就应只涉及阴影，不应让模型重画主体。使用时可以把模型自己写出的偏差说明与要求记录逐项对照：它漏掉的偏差不会被修正，它误报的偏差可能引入新的错误。

修订预算需要事先规定。UniReason原论文的实现只作一轮修订，不能外推为“修订越多次越好”。每一轮都要额外付出一次图像编码、一段文字生成、一次完整图像采样和解码，计算量大约是一次直接生成的两倍以上。更重要的是，修订可能在纠正数量的同时引起主体漂移，例如狗的斑纹改变。所以合理的规则是设定修订轮数上限（例如一轮），并保存原草图；若修订结果在已满足的项上退步，就拒绝修订、退回草图，或者转为下一节的局部编辑。

验收应固定原条件，把直接生成、加入文字推理和再加一轮修订的三个结果并排比较，逐项检查数量、颜色归属、左右关系和阴影，然后才看自然度。若偏差说明全部正确而修订图仍画错关系，应记为图像实施失败，而不能因文字推理正确就算任务完成。UniReason训练时用外部大模型核验草图并构造修订数据，验收时若仍用同类模型裁判打分，结果可能只是更迎合这个裁判；因此还要加入区域检测、字符读取或人工抽查等独立检查。

\section{已有图像的转换}
\label{sec:vision-generation-translation}

转换以一幅已有图像为输入，输出另一种风格、呈现或成像形式的图像，例如把玩具狗照片转成水彩画、把边缘草图转成照片、把白天街景转成夜景。与创建相比，输入图像带来了一个新的约束：输出必须与输入保持某种对应，例如狗的位置、姿态和数量不变；与下一节的编辑相比，转换作用于整幅图像，变化的方式由事先确定的“域”决定，而不是由每次的请求临时规定。难点在于怎样学到一个既改变外观、又保持对应的映射，而能够学到什么样的对应，取决于训练数据提供了什么。

\subsection{转换要求的规定}
\label{sec:vision-generation-translation-requirements}

转换要求要写清两件事：输出属于什么目标形式，以及哪些内容必须与输入对应。把玩具狗照片转成水彩画时，纹理、笔触和色调可以改变，狗的姿态、数量和画面布局要保持；把边缘图转成照片时，每条边缘的位置都必须保留，颜色和材质可以由模型补充；若把一种成像方式转换成另一种，例如把卫星图转成地图，还要求每个位置的结构与实际地物对应。三种用途即使用相同网络实现，允许的误差也不同。

训练数据必须与这种要求相符。对齐的“边缘图—照片”配对，逐像素地告诉模型输入的每个位置应该变成什么；只有一组照片和一组水彩画、二者并不对应时，数据只说明两个域各自的外观；只有一幅指定风格的画时，数据提供的是纹理统计，而不是任何一幅图的目标。若把没有逐位置对应的图像当作像素监督，损失就会迫使网络拟合本不存在的对应关系。因此，转换要求、可用数据和下面的模型选择是联动的。

\subsection{转换映射的学习}
\label{sec:vision-generation-translation-mapping}

转换映射是一个前馈网络$G$，读入一幅图像，一次前向输出转换结果。三种模型的差别在于用什么监督来约束“输出与输入的对应”。pix2pix有逐像素配对的目标图像，直接用目标图约束输出，并用判别器约束局部外观；CycleGAN没有配对，只能要求转换后再转回来能恢复输入，用这种往返关系代替配对目标；Johnson等人的感知损失网络面对的是固定风格转换，既没有配对目标也不需要另一个域的数据集，它用一个固定的预训练网络在特征空间中分别度量“内容是否保留”和“风格是否接近”。这三条路线依次放宽了对训练数据的要求，也依次减弱了对输出的约束。

\subsubsection{pix2pix}
\label{sec:vision-generation-pix2pix}

卷五的“基于神经网络的生成模型”章已经完整讲了pix2pix的结构与目标：它是条件生成对抗网络，生成器是带跳连的U-Net，判别器是读入“输入—输出”对的局部块判别器，目标是条件对抗项加上$\lambda$倍的$\ell_1$像素误差\scite{vision-pix2pix2017}。这里回顾它要解决的问题，然后侧重在转换任务中怎样组织数据、选择超参数和使用。

pix2pix的动机是：对于有配对数据的转换，只用$\ell_1$像素误差训练，网络会对“同一输入可能对应的多种纹理”取平均，输出模糊；只用对抗项，输出逼真但可能不保留输入结构。它的主要想法是两者兼用：$\ell_1$项把输出拉向配对目标，保证低频结构对应；判别器只看局部图像块，专门惩罚模糊和不真实的纹理。判别器同时读入输入，使它判断的是“这个输出是否像这个输入的真实对应”，而不仅是“像不像照片”。

以$256\times256$的边缘图转玩具狗照片为例，逐层结构见表\ref{tab:vision-pix2pix-structure}。论文附录用两个记号描述各层：C$k$表示“卷积—批归一化—ReLU”，有$k$个卷积核；CD$k$在批归一化之后多一个丢弃率50\%的dropout。所有卷积都是$4\times4$，编码器和判别器中步长为2、每层把分辨率减半，解码器中每层把分辨率加倍。编码器的ReLU都是斜率0.2的leaky ReLU，第一层C64不加批归一化；解码器使用普通ReLU。

\begin{table*}[tbp]
 \centering
 \small
 \begin{tabularx}{\linewidth}{@{}p{30mm}Xp{24mm}p{22mm}p{12mm}p{20mm}@{}}
  \toprule
  阶段或层组 & 层类型 & 卷积核与步长 & 输出通道 & 重复次数 & 输出尺寸\\
  \midrule
  \multicolumn{6}{@{}l}{\textit{生成器（U-Net），输入$256\times256\times3$}}\\
  编码层1 & 卷积，leaky ReLU（无批归一化） & $4\times4$，步长2 & 64 & 1 & $128\times128$\\
  编码层2 & C128 & $4\times4$，步长2 & 128 & 1 & $64\times64$\\
  编码层3 & C256 & $4\times4$，步长2 & 256 & 1 & $32\times32$\\
  编码层4至8 & C512 & $4\times4$，步长2 & 512 & 5 & $16\times16\to1\times1$\\
  解码层1至3 & 上采样卷积，批归一化，dropout，ReLU，再拼接同分辨率编码特征 & $4\times4$，步长$1/2$ & 512，拼接后1024 & 3 & $2\times2\to8\times8$\\
  解码层4 & 同上，无dropout & $4\times4$，步长$1/2$ & 512，拼接后1024 & 1 & $16\times16$\\
  解码层5至7 & 同上，无dropout & $4\times4$，步长$1/2$ & 256、128、64，拼接后512、256、128 & 3 & $32\times32\to128\times128$\\
  输出层 & 上采样卷积，tanh & $4\times4$，步长$1/2$ & 3 & 1 & $256\times256$\\
  \midrule
  \multicolumn{6}{@{}l}{\textit{$70\times70$判别器，输入为输入图与输出图沿通道拼接的$256\times256\times6$}}\\
  C64 & 卷积，leaky ReLU（无批归一化） & $4\times4$，步长2 & 64 & 1 & $128\times128$\\
  C128 & 卷积，批归一化，leaky ReLU & $4\times4$，步长2 & 128 & 1 & $64\times64$\\
  C256 & 卷积，批归一化，leaky ReLU & $4\times4$，步长2 & 256 & 1 & $32\times32$\\
  C512 & 卷积，批归一化，leaky ReLU & $4\times4$，步长1（代码） & 512 & 1 & $31\times31$\\
  输出层 & 卷积，sigmoid & $4\times4$，步长1（代码） & 1 & 1 & $30\times30$\\
  \bottomrule
 \end{tabularx}
 \caption{pix2pix的生成器与$70\times70$判别器。层序列、通道与激活取自论文附录：编码器C64-C128-C256-C512-C512-C512-C512-C512，U-Net解码器按论文记号为CD512-CD1024-CD1024-C1024-C1024-C512-C256-C128，记号中的数字是各解码层的输入通道数，即上一层输出与跳连拼接之后的通道数，最后的C128对应输出层的输入。论文写所有卷积步长为2，但那样判别器的感受野不是$70\times70$；表中判别器最后两层的步长1与$30\times30$输出取自作者公开的代码。}
 \label{tab:vision-pix2pix-structure}
\end{table*}

数据流如下。输入$\bm{X}\in\mathbb{R}^{256\times256\times3}$经编码器8层步长为2的卷积，逐步缩小到$1\times1\times512$；解码器对称地8次上采样，第$i$个解码层的输出与第$8-i$个编码层的输出沿通道拼接后送入下一层，所以论文把解码层的通道写成1024、512等拼接后的数目。最终输出$\widehat{\bm{Y}}\in\mathbb{R}^{256\times256\times3}$，经tanh落在$[-1,1]$。跳连使边缘位置可以直接从编码器传到解码器，而不必经过$1\times1$的瓶颈；没有跳连的编码—解码版本，解码器为CD512-CD512-CD512-C512-C256-C128-C64，所有信息都必须穿过瓶颈。判别器把$\bm{X}$与$\widehat{\bm{Y}}$（或真实$\bm{Y}$）沿通道拼成6通道输入，输出一个$30\times30$的分数图。每个分数的感受野可以逐层累加：对$4\times4$卷积，第$\ell$层使感受野增加$3\times(\text{前面各层步长之积})$，五层的步长依次为$2,2,2,1,1$，感受野为$1+3\times(1+2+4+8+8)=70$。所以每个分数只判断输入上一个$70\times70$邻域内的配对是否真实。论文还比较了感受野为$1\times1$、$16\times16$和$286\times286$的判别器，它们只是增减了C层的数目，例如$16\times16$判别器为C64-C128。

训练时交替更新：判别器以真实对为正、生成对为负，在分数图的每个位置上计算对数损失再平均；生成器使用非饱和对抗项，并加上$\lambda\lVert\bm{Y}-\widehat{\bm{Y}}\rVert_1$，原论文取$\lambda=100$。优化判别器时把目标除以2，减慢判别器相对生成器的学习；两者都用Adam，学习率$2\times10^{-4}$，$\beta_1=0.5,\beta_2=0.999$，权重从均值0、标准差0.02的高斯分布初始化，全部从头训练。数据增强先把$256\times256$的图像放大到$286\times286$，再随机裁回$256\times256$。可以手算一次：设两个局部块的真实分数为$(0.8,0.6)$、生成分数为$(0.2,0.4)$，判别损失为$-\{\log0.8+\log0.6+\log0.8+\log0.6\}/2\approx0.734$（自然对数）；生成器的非饱和对抗项为$-(\log0.2+\log0.4)/2\approx1.263$。像素项是求和还是求平均必须与$\lambda$配套说明，否则权重不可比较。数据增强必须对输入与目标做相同的裁剪和翻转，保证配对位置不变。

使用时只运行生成器，判别器和目标图都不再需要。原实现在测试时保留dropout，并用当前测试批次计算批归一化统计（批大小为1时相当于实例归一化），复现时要保存这些设置。常见失败有三种：训练配对存在系统性错位时，网络会学会这种错位；同一输入有多种合理纹理时，输出缺乏多样性，原论文发现加入随机噪声会被网络忽略，dropout带来的变化也很有限；局部块逼真不保证整幅图的结构和身份正确。

\subsubsection{CycleGAN}
\label{sec:vision-generation-cyclegan}

若要把玩具狗照片转成水彩画，几乎不可能找到同一只狗、同一姿态的“照片—水彩”配对，pix2pix的$\ell_1$项就无从计算。CycleGAN处理的正是只有两组互不对应的图像的情形\scite{vision-cyclegan2017}；其目标和推导同样已在卷五该章讲过。它的主要想法是：同时训练照片到水彩的$G:A\to B$和水彩到照片的$F:B\to A$，用目标域判别器要求转换结果像目标域图像，再用\term{循环一致性}（cycle consistency）要求$F(G(\bm{x}))\approx\bm{x}$、$G(F(\bm{y}))\approx\bm{y}$。往返能恢复输入，说明转换没有丢掉输入的信息，这就替代了缺失的配对目标。

循环损失为
\begin{equation}
 L_{\mathrm{cyc}}=
 \mathbb{E}_{\bm{x}\sim p_A}
 [\lVert F(G(\bm{x}))-\bm{x}\rVert_1]
 +
 \mathbb{E}_{\bm{y}\sim p_B}
 [\lVert G(F(\bm{y}))-\bm{y}\rVert_1]\text{。}
 \label{eq:vision-cycle-consistency}
\end{equation}
总目标是两个方向的对抗项加上$\lambda L_{\mathrm{cyc}}$，原论文取$\lambda=10$。架构上，两个生成器沿用Johnson等人的转换网络（见下一小节），逐层结构见表\ref{tab:vision-cyclegan-structure}。论文附录沿用Johnson代码库的命名：c7s1-$k$是$7\times7$、步长1、$k$个卷积核的“卷积—实例归一化—ReLU”；d$k$是$3\times3$、步长2的同类层，用于下采样；R$k$是含两层$3\times3$卷积、两层卷积核数相同的残差块；u$k$是$3\times3$、步长$1/2$的分数步长卷积层，用于上采样。卷积使用反射填充以减少边界伪影。对$256\times256$及更高分辨率的训练图像，生成器为c7s1-64、d128、d256、9个R256、u128、u64、c7s1-3；对$128\times128$的图像，残差块减为6个。实例归一化对每幅图、每个通道单独计算均值和方差，与批大小1的训练相配。两个判别器都是$70\times70$的局部块判别器，层序列为C64-C128-C256-C512，C$k$为$4\times4$的“卷积—实例归一化—leaky ReLU（斜率0.2）”，第一层不加实例归一化，最后再用一层卷积输出单通道分数图。与pix2pix不同，它不读入输入图像，只判断一幅图是否属于本域。

\begin{table*}[tbp]
 \centering
 \small
 \begin{tabularx}{\linewidth}{@{}p{30mm}Xp{24mm}p{22mm}p{12mm}p{20mm}@{}}
  \toprule
  阶段或层组 & 层类型 & 卷积核与步长 & 输出通道 & 重复次数 & 输出尺寸\\
  \midrule
  \multicolumn{6}{@{}l}{\textit{生成器（9个残差块），输入$256\times256\times3$}}\\
  c7s1-64 & 卷积，实例归一化，ReLU & $7\times7$，步长1 & 64 & 1 & $256\times256$\\
  d128 & 卷积，实例归一化，ReLU & $3\times3$，步长2 & 128 & 1 & $128\times128$\\
  d256 & 卷积，实例归一化，ReLU & $3\times3$，步长2 & 256 & 1 & $64\times64$\\
  R256 & 残差块，两层卷积 & $3\times3$，步长1 & 256 & 9 & $64\times64$\\
  u128 & 分数步长卷积，实例归一化，ReLU & $3\times3$，步长$1/2$ & 128 & 1 & $128\times128$\\
  u64 & 分数步长卷积，实例归一化，ReLU & $3\times3$，步长$1/2$ & 64 & 1 & $256\times256$\\
  c7s1-3 & 卷积 & $7\times7$，步长1 & 3 & 1 & $256\times256$\\
  \midrule
  \multicolumn{6}{@{}l}{\textit{$70\times70$判别器，输入$256\times256\times3$}}\\
  C64至C512 & 卷积，实例归一化（C64除外），leaky ReLU & $4\times4$，见表\ref{tab:vision-pix2pix-structure} & 64、128、256、512 & 4 & $128\times128\to31\times31$\\
  输出层 & 卷积 & $4\times4$，步长1（代码） & 1 & 1 & $30\times30$\\
  \bottomrule
 \end{tabularx}
 \caption{CycleGAN的生成器与判别器，层序列取自论文附录。判别器与pix2pix的$70\times70$判别器结构相同，只是把批归一化换成实例归一化、输入只有3通道；各层步长与输出尺寸的说明同表\ref{tab:vision-pix2pix-structure}。附录对最后一层c7s1-3的输出激活没有单独说明。}
 \label{tab:vision-cyclegan-structure}
\end{table*}

论文只说明残差块含两层$3\times3$卷积；按作者公开的实现，单个残差块内的计算顺序是：$3\times3$卷积、实例归一化、ReLU、$3\times3$卷积、实例归一化，再与块的输入相加。9个残差块都在$64\times64\times256$上工作，转换主要在这个分辨率上完成；下采样把感受野扩大，上采样再恢复到原分辨率。由于整个网络都是卷积，训练后可以处理与训练尺寸不同的输入。

训练时从两个域独立抽样，批大小为1。对抗项采用最小二乘形式：以域$B$为例，判别器最小化$\frac12\mathbb{E}[(D_B(\bm{y})-1)^2]+\frac12\mathbb{E}[D_B(G(\bm{x}))^2]$，生成器最小化$\mathbb{E}[(D_B(G(\bm{x}))-1)^2]$，平方误差在判别器很自信时仍能给生成器提供梯度。判别器用最近生成的50幅图像组成的缓冲池更新，减弱它只追随最新生成器的振荡。原实现用Adam、学习率$2\times10^{-4}$训练100轮，再在接下来100轮线性衰减到零。对照片转绘画这类任务，原论文还加入“输入已属于目标域时尽量不改动”的恒等映射项，以保持颜色构成。使用时只运行需要的方向$G$，反向生成器和两个判别器只在训练中使用。

循环一致性有一个根本限制：它约束信息不丢失，却不唯一确定语义对应。设两个域各只有一幅狗图和一幅猫图，概率各为$1/2$。若$G$把狗映成猫、猫映成狗，$F$执行逆映射，那么两个输出分布都正确、循环误差为零，但语义完全交换。图\ref{fig:vision-cycle-ambiguity}把这个反例与保持类别的映射并排显示。实际训练中，卷积网络的局部性倾向于保持布局，这类交换较少出现，但这是网络结构带来的倾向，不是目标函数的保证。

\begin{figure}[htbp]
 \centering
 % Both maps are bijections of two equiprobable categories.
% Target order is changed in the right panel to avoid line crossings.
\begin{tikzpicture}[x=1mm,y=1mm,font=\figurefont\small,
 b/.style={draw=BookBlue,fill=BookBlue!8,line width=.8pt,
  minimum width=12mm,minimum height=8mm,inner sep=1mm},
 a/.style={{Stealth[length=1.7mm]}-{Stealth[length=1.7mm]},
  BookInk,line width=.9pt}]
 \node[font=\figurefont\small\bfseries] at (20,14) {保持类别};
 \node[font=\figurefont\small\bfseries] at (79,14) {交换类别};
 \node[anchor=south,text=BookMuted] at (5,5) {域$A$};
 \node[anchor=south,text=BookMuted] at (35,5) {域$B$};
 \node[anchor=south,text=BookMuted] at (64,5) {域$A$};
 \node[anchor=south,text=BookMuted] at (94,5) {域$B$};
 \node[b] at (5,0) (ad) {狗};
 \node[b] at (35,0) (bd) {狗};
 \node[b] at (5,-17) (ac) {猫};
 \node[b] at (35,-17) (bc) {猫};
 \node[b] at (64,0) (ad2) {狗};
 \node[b] at (94,0) (bc2) {猫};
 \node[b] at (64,-17) (ac2) {猫};
 \node[b] at (94,-17) (bd2) {狗};
 \foreach \u/\v in {ad/bd,ac/bc,ad2/bc2,ac2/bd2} {
  \draw[a] (\u)--node[above=1mm] {$G/F$}(\v);
 }
 \node[anchor=north] at (20,-25) {循环误差为0};
 \node[anchor=north] at (79,-25) {循环误差也为0};
\end{tikzpicture}

 \caption{循环一致性不能唯一确定语义对应。两域各有等概率的狗、猫两项。左侧保持类别，右侧交换类别；每条往返路径都回到原项，且目标域边缘分布相同。反例是离散教学构造，用来说明循环与域判别约束仍有多解。}
 \label{fig:vision-cycle-ambiguity}
\end{figure}

CycleGAN的常见失败也由此而来。转换需要大幅改变几何形状时，例如把猫转成狗，局部纹理改变容易而形状改变困难，结果往往只改了纹理；生成器还可能把输入信息以人眼不可见的细微信号藏在输出中，使往返误差很小而中间图像并不忠实。因此在测量或医学等需要真实对应的转换中，无配对训练不能提供保真保证。

\subsubsection{Johnson感知损失网络}
\label{sec:vision-generation-perceptual-style}

第三种情形是：只有一幅指定的水彩画，希望任何输入照片都能转成它的风格。Gatys等人的风格迁移已经表明，在预训练卷积网络的特征空间中，用特征本身度量内容、用特征的通道相关统计度量风格，就能把一幅图的内容与另一幅图的风格组合起来\scite{vision-gatys2016}；但那种方法对每幅新照片都要从噪声出发、用数百次梯度迭代直接优化像素，每幅图都要付出一次完整的优化。Johnson等人的动机是保留这种特征空间的度量，却把逐图优化改成一次前向计算\scite{vision-perceptual2016}。

主要想法是区分两个网络。\term{感知损失}（perceptual loss）指在固定预训练网络的特征空间中计算的差异，而不是像素差异。可训练的图像变换网络$G_{\bm{\theta}}$负责生成图像；固定的损失网络（ImageNet上预训练的VGG-16）只用于计算感知损失，不更新参数。用许多内容图训练$G_{\bm{\theta}}$，使它的输出在损失网络看来内容接近输入、风格接近指定画作；训练完成后，转换一幅新图只需$G_{\bm{\theta}}$的一次前向，原论文报告与逐图优化的质量相当，而速度快约三个数量级。

风格转换用的变换网络逐层结构见表\ref{tab:vision-johnson-structure}，取自论文的补充材料。输入$256\times256\times3$的照片，先经一层$9\times9$卷积得到32通道，再经两次步长为2的$3\times3$卷积下采样到$64\times64\times128$；接着5个残差块在这个分辨率上处理；然后两次步长为$1/2$的$3\times3$卷积恢复到$256\times256$，最后一层$9\times9$卷积输出3通道。除输出层外，每个非残差卷积层后面都接空间批归一化和ReLU；输出层改用缩放的tanh，把像素限制在$[0,255]$。网络没有池化层，下采样和上采样都由带步长的卷积完成。全部为卷积层，所以训练后可以处理任意尺寸的输入。

\begin{table}[htbp]
 \centering
 \small
 \begin{tabularx}{\linewidth}{@{}p{16mm}Xp{15mm}p{9mm}p{6mm}p{15mm}@{}}
  \toprule
  阶段或层组 & 层类型 & 卷积核与步长 & 输出通道 & 重复次数 & 输出尺寸\\
  \midrule
  输入层 & 卷积，批归一化，ReLU & $9\times9$，步长1 & 32 & 1 & $256\times256$\\
  下采样 & 卷积，批归一化，ReLU & $3\times3$，步长2 & 64、128 & 2 & $128\times128\to64\times64$\\
  残差块 & 两层卷积，相加后不接ReLU & $3\times3$，步长1 & 128 & 5 & $64\times64$\\
  上采样 & 分数步长卷积，批归一化，ReLU & $3\times3$，步长$1/2$ & 64、32 & 2 & $128\times128\to256\times256$\\
  输出层 & 卷积，缩放tanh & $9\times9$，步长1 & 3 & 1 & $256\times256$\\
  \bottomrule
 \end{tabularx}
 \caption{Johnson等人用于风格转换的变换网络，取自论文补充材料。按层数计共16层（5个残差块各含2层）。同一论文用于超分辨率的网络结构不同：没有下采样，残差块为64通道，再用2或3次步长$1/2$的卷积放大4倍或8倍。}
 \label{tab:vision-johnson-structure}
\end{table}

残差块采用Gross与Wilber的设计：$3\times3$卷积、批归一化、ReLU、$3\times3$卷积、批归一化，与输入相加，相加之后不再接ReLU。补充材料对比了把每个残差块换成同样两层卷积的非残差网络，二者最终的训练损失相近，残差网络收敛更快；作者推测残差连接对更深的网络更关键，这个16层的网络较浅。

损失由三项组成。设损失网络第$\ell$层的特征为$\bm{F}_\ell(\bm{X})\in\mathbb{R}^{c_\ell\times n_\ell}$，$c_\ell$为通道数，$n_\ell$为空间位置数。内容项比较输出$\widehat{\bm{X}}$与输入$\bm{X}$在relu2\_2层的特征：$\frac{1}{c_\ell n_\ell}\lVert\bm{F}_\ell(\widehat{\bm{X}})-\bm{F}_\ell(\bm{X})\rVert_{\mathrm{F}}^2$。风格项比较输出与风格图$\bm{S}$的Gram矩阵
\begin{equation}
 \bm{\Gamma}_\ell(\bm{X})
 =\frac{1}{c_\ell n_\ell}
   \bm{F}_\ell(\bm{X})\bm{F}_\ell(\bm{X})^\top\text{，}
 \label{eq:vision-style-gram}
\end{equation}
在relu1\_2、relu2\_2、relu3\_3、relu4\_3四层上累加$\lVert\bm{\Gamma}_\ell(\widehat{\bm{X}})-\bm{\Gamma}_\ell(\bm{S})\rVert_{\mathrm{F}}^2$。第三项是总变差正则，惩罚输出相邻像素的过大跳变，使图像平滑。三项加权求和，权重决定内容与风格的折中。

为什么内容项用特征本身、风格项用Gram矩阵？Gram矩阵对空间位置求和，只记录哪些通道同时激活，不记录它们出现在哪里。用一个两通道、两位置的例子验证：内容特征$\bm{F}=\left(\begin{smallmatrix}1&0\\0&2\end{smallmatrix}\right)$给出$\bm{\Gamma}=\operatorname{diag}(1/4,1)$；把两个位置对调得到$\bm{F}'=\left(\begin{smallmatrix}0&1\\2&0\end{smallmatrix}\right)$，Gram矩阵不变，但内容特征的平方差和为10，按$c_\ell n_\ell=4$归一化为2.5。所以风格项只要求纹理统计相近，允许笔触出现在任何位置；内容项才把结构锁定在输入的相应位置。

训练数据是约8万幅MS-COCO图像，缩放为$256\times256$，与一幅固定风格图配合。每个小批量（大小为4）只抽内容图，经$G_{\bm{\theta}}$得到输出，再让损失网络分别处理输出与内容图，风格图的Gram矩阵可预先计算。梯度经过损失网络反传到输出图像，再传到$G_{\bm{\theta}}$，但损失网络本身不更新。原论文用Adam、学习率$10^{-3}$训练4万次迭代，约两轮；总变差正则的权重在$10^{-6}$到$10^{-4}$之间按每种风格交叉验证选取，不用权重衰减和dropout，单张GTX Titan X上训练约4小时。这种网络一个模型只对应一种风格，换一幅风格画要重新训练；它也不保证徽标、细小文字等精确结构的保留。

三种转换模型的比较见表\ref{tab:vision-generation-translation-comparison}。固定任务为“把玩具狗照片转成某种绘画”，并以“姿态与布局是否保持、是否像目标形式”为评价口径。CycleGAN相对pix2pix改变的是监督来源：用往返重建替代配对目标，换来对无配对数据的适用性，代价是语义对应不再唯一，还要多训练一个反向生成器和一个判别器。Johnson网络又把“另一个域的数据集”换成“一幅风格图加一个固定特征网络”，彻底去掉了判别器和对抗训练，训练稳定，但只能表达纹理统计层面的风格，不能完成需要改变语义或几何的转换。

\begin{table}[htbp]
 \centering
 \small
 \begin{tabularx}{\linewidth}{@{}l X X X@{}}
  \toprule
   & pix2pix & CycleGAN & Johnson感知损失网络\\
  \midrule
  训练数据 & 逐像素对齐的输入—目标配对 & 两个域互不对应的图像集 & 内容图集与一幅风格图\\
  对应关系的来源 & 配对目标的$\ell_1$项 & 往返重建的循环项 & 固定网络的内容特征项\\
  目标外观的来源 & 条件局部块判别器 & 目标域局部块判别器 & Gram矩阵风格项\\
  训练时的网络 & 生成器、判别器 & 两个生成器、两个判别器 & 变换网络、固定损失网络\\
  使用时的网络 & 生成器 & 一个方向的生成器 & 变换网络\\
  主要局限 & 依赖对齐配对，输出缺乏多样性 & 语义对应不唯一，难改几何形状 & 一个模型一种风格，只改纹理统计\\
  \bottomrule
 \end{tabularx}
 \caption{三种转换映射在同一转换任务下的比较。三者使用时都是一次前向计算，差别主要在训练数据与监督。}
 \label{tab:vision-generation-translation-comparison}
\end{table}

\subsection{转换结果的核验}
\label{sec:vision-generation-translation-check}

转换结果要同时满足两个方面：具有目标形式，并保持规定的对应。输出仍像照片，说明风格没有充分改变；输出是自然的水彩画但狗变成了猫，说明语义保持失败。两方面各需独立检查，一方的好结果不能抵消另一方的失败。CycleGAN的循环误差很小，只回答了往返约束，不能代替对中间结果的检查；Johnson网络的感知损失很低，也只说明在VGG特征下接近，不说明徽标和文字保留。

具体做法是固定一组未参与训练的输入，逐项记录对象、轮廓、文字与颜色的变化，按用途区分允许和不允许的改变。轮廓保持可以用边缘图或分割掩码在输入与输出之间比较；对象身份可以用\ref{chap:vision-image-recognition}章的检测器检查类别与数量；目标形式可以用目标域的分类器或人工判断。若发现系统性错位，应先检查训练配对；若没有配对依据却要求严格的成像保真，就需要增加相应监督，或者收窄用途。转换由预先确定的域关系决定变化方式；若每次请求要求改变的内容不同，就进入下一节的编辑。

\section{指定内容的编辑}
\label{sec:vision-generation-editing}

编辑的输入是一幅原图和一个请求，例如“把狗旁边的杯子换成花瓶”，输出是只在请求涉及的部分发生变化的图像。它与转换的区别在于变化由当前请求规定，而且通常只覆盖局部；与创建的区别在于区域外的内容是必须保持的约束。难点有三处：确定改哪里，在改动区域内生成与周围协调的新内容同时保持其余部分，以及在多轮请求中不丢失已经接受的修改。本节按编辑流程的步骤展开，最后讨论核验。

\subsection{编辑范围的确定}
\label{sec:vision-generation-editing-scope}

“把玩具狗旁边的杯子换成花瓶”至少包含三个要素：待替换的对象、替换后的内容和应保留的周围物体。若画面里有两个杯子，仅凭类别词无法确定改哪一个。用户可以直接给出框、掩码或参考编号，也可以从\ref{chap:vision-image-recognition}章的检测与分割结果中选择；对“左边那个蓝色杯子”这类指代的自动解析由\ref{chap:vision-multimodal-understanding}章展开，不是编辑算法的前提。

范围通常以掩码表示。沿用\ref{chap:vision-image-restoration}章的约定，$\bm{M}\in\{0,1\}^{h\times w}$中$m_{i,j}=1$表示保留区，$m_{i,j}=0$表示允许生成区。若花瓶比原杯子高，只遮住杯子轮廓就没有容纳新物体的空间，还会留下杯子原来的阴影；合理的范围应覆盖新物体可能占据的位置、接触区和阴影，并说明玩具狗是否允许被遮挡。掩码是任务要求的一部分，生成器不会自动推断出正确的编辑边界。

文字也可能不足以说明形状或位置。视觉指令把箭头、圈选、骨架或草图作为请求的一部分：圈住杯子并写“换成花瓶”时，圈选用于选择目标；画出一个更高的花瓶轮廓时，草图规定形状；画一个侧向箭头表示光线时，要求的是相应的受光与阴影，而箭头本身不应出现在结果中。VIBE基准按指示定位、形态修改和因果关系三层组织这类编辑任务，用多模态模型裁判按任务专用问题判分\scite{vision-vibe-2026}。它是评价材料，不是编辑模型。用它检验时，可以让文字相同而箭头分别指向左右两个杯子，检查被替换的对象是否随之改变；若只看整图是否出现花瓶，两次请求可能都得高分，视觉指令是否起作用却没有被检验。

位置标注必须跟随前处理变换。原图为$1000\times800$像素，目标点为$(600,400)$；若等比缩小一半，点变为$(300,200)$；若再从左侧裁去50像素，模型输入中的点是$(250,200)$。输出恢复时要逆向补回裁剪偏移并还原比例，掩码与草图也作相同变换。只缩放图像而保留原坐标，会把编辑送到错误位置。

\subsection{单轮编辑的实施}
\label{sec:vision-generation-editing-methods}

有了原图、请求和范围，单轮编辑要回答两个问题：新内容从哪里来，以及怎样让它与周围衔接。经典路线只从已有图像中搬运内容。泊松图像编辑把另一幅照片中现成的花瓶贴进范围，在梯度域中求解，使边界处没有接缝；Criminisi样例修补用同一幅图中别处的图像块，填补移除杯子后留下的空洞，并按优先级决定先填哪里；PatchMatch解决这两类块方法共同的计算瓶颈，即为每个图像块快速找到最相似的块。学习路线按文字生成原图中不存在的内容。GLIDE修补模型要求用户给出掩码，只在掩码内生成与文字和周围协调的内容；InstructPix2Pix只要一句指令，由模型自己判断哪里该改、改成什么。沿着这个顺序，新内容的来源从“现成图像”变成“生成模型”，范围的给出者从用户变成模型：经典路线对保留区的约束最严格，却只能重组已有内容；学习路线能创造新内容，但需要大量训练数据，指令路线还可能改动未指定的区域。

\subsubsection{泊松图像编辑}
\label{sec:vision-generation-poisson}

假设手头有一张现成的花瓶照片，想把花瓶放到玩具狗照片里杯子的位置。直接把花瓶区域的像素复制过去，两张照片的曝光和色温不同，花瓶边缘会出现一圈明显的接缝；在边界附近用渐变的权重混合两幅图，接缝变得模糊，却会留下一圈半透明的光晕，花瓶自身的细节也被冲淡。Pérez等人的出发点是，人眼对局部的明暗变化比对整体亮度更敏感：只要花瓶内部的明暗起伏保留下来，整体亮度略有偏移并不显眼\scite{vision-poisson2003}。于是问题改写为：区域内的结果应尽量保持源图像的梯度，区域边界上的值必须与目标图像相等。

记目标图像（玩具狗照片）为$f^*$，源图像（花瓶照片，已对齐到目标图像的坐标）为$g$，要粘贴的区域为$\Omega$，其边界为$\partial\Omega$；未知的是$\Omega$内的结果$f$。各颜色通道分别处理，下面只写一个通道。给定区域内的一个向量场$\bm{v}$作为引导，称为引导场（guidance field），泊松图像编辑求
\begin{equation}
 \min_{f}\iint_{\Omega}\lVert\nabla f-\bm{v}\rVert_2^2\,\mathrm{d}x\,\mathrm{d}y,
 \qquad f|_{\partial\Omega}=f^*|_{\partial\Omega}\text{。}
 \label{eq:vision-poisson-variational}
\end{equation}
把花瓶无缝地贴进来，取$\bm{v}=\nabla g$，这种用法称为无缝克隆。$\bm{v}=\bm{0}$时，问题变成在固定边界值下求最平缓的填充，解满足拉普拉斯方程$\Delta f=0$，其中$\Delta f=\partial^2f/\partial x^2+\partial^2f/\partial y^2$。这正是Dirichlet能量最小化：在边界值固定的函数中，使$\iint|\nabla f|^2$最小的函数满足拉普拉斯方程，可参考\BookChapterRef{chap:functional-analysis-and-approximation}{卷一“函数空间”章}。加入引导场以后，推导步骤相同，只多出一个源项。

设$f$是最小解，$\eta$是$\Omega$上任意一个光滑且在$\partial\Omega$上为0的函数。对任意实数$\epsilon$，$f+\epsilon\eta$仍满足边界条件，因此$J(\epsilon)=\iint_{\Omega}\lVert\nabla f+\epsilon\nabla\eta-\bm{v}\rVert_2^2$在$\epsilon=0$处取最小，$J'(0)=0$。展开平方并对$\epsilon$求导得
\begin{equation*}
 J'(0)=2\iint_{\Omega}(\nabla f-\bm{v})\cdot\nabla\eta\,\mathrm{d}x\,\mathrm{d}y=0\text{。}
\end{equation*}
对向量场$\bm{u}=\nabla f-\bm{v}$，乘积法则给出$\operatorname{div}(\eta\bm{u})=\nabla\eta\cdot\bm{u}+\eta\operatorname{div}\bm{u}$；对左边用散度定理，区域积分化为边界上的通量$\oint_{\partial\Omega}\eta\,\bm{u}\cdot\bm{n}\,\mathrm{d}s$，其中$\bm{n}$是边界的外法向，而$\eta$在边界上为0，这一项为0。于是
\begin{equation*}
 \iint_{\Omega}\nabla\eta\cdot(\nabla f-\bm{v})
 =-\iint_{\Omega}\eta\,\bigl(\Delta f-\operatorname{div}\bm{v}\bigr)=0
\end{equation*}
对所有这样的$\eta$成立。若连续函数$\Delta f-\operatorname{div}\bm{v}$在某点不为0，取一个只在该点附近非零、与它同号的$\eta$，积分就不为0，矛盾；所以最小解满足带源项的泊松方程
\begin{equation}
 \Delta f=\operatorname{div}\bm{v}\quad\text{在}\ \Omega\ \text{内},
 \qquad f|_{\partial\Omega}=f^*|_{\partial\Omega}\text{。}
 \label{eq:vision-poisson-equation}
\end{equation}
无缝克隆时$\operatorname{div}\bm{v}=\Delta g$，记修正量$\tilde f=f-g$，方程变为$\Delta\tilde f=0$，边界值为$f^*-g$。这给出了一个直观的解释：结果等于源图像加上一层平缓的修正，修正量在边界上恰好抹平两幅图的差，再向内部平滑过渡。花瓶内部的细节由$g$保留，整体颜色则被边界上的差带偏，这既是接缝消失的原因，也是被贴物体颜色会改变的原因。这里用到的多元变分推导与带源项的泊松方程，前两卷只讲了无源项的特例，本章给出的推导已足够当前使用。

实际计算在像素网格上进行。设$\Omega$是一组像素，$N_p$是像素$p$的4个邻居，$\partial\Omega$是不属于$\Omega$但与$\Omega$相邻的像素。对每对相邻像素$(p,q)$规定引导量$v_{p,q}$，满足$v_{q,p}=-v_{p,q}$，无缝克隆取$v_{p,q}=g_p-g_q$。离散问题是
\begin{equation*}
 \min_{\{f_p\}_{p\in\Omega}}
 \sum_{\langle p,q\rangle\cap\Omega\neq\varnothing}
 (f_p-f_q-v_{p,q})^2,
 \qquad f_q=f^*_q\ \ (q\in\partial\Omega)\text{，}
\end{equation*}
求和遍历至少有一端在$\Omega$内的相邻像素对，每对只算一次。对$p\in\Omega$求偏导：含$f_p$的项正好是$p$与它的各个邻居组成的对，每一项$(f_p-f_q-v_{p,q})^2$对$f_p$的导数是$2(f_p-f_q-v_{p,q})$（若这一对在求和中写成$(q,p)$，由$v_{q,p}=-v_{p,q}$得到同一结果）。令偏导为0，把已知的边界值移到右边，得到
\begin{equation}
 |N_p|\,f_p-\sum_{q\in N_p\cap\Omega}f_q
 =\sum_{q\in N_p\cap\partial\Omega}f^*_q
 +\sum_{q\in N_p}v_{p,q},
 \qquad p\in\Omega\text{。}
 \label{eq:vision-poisson-discrete}
\end{equation}
左边是离散拉普拉斯算子的相反数：内部像素$|N_p|=4$，即“4倍中心值减去四邻之和”，图像锐化中用到的就是这个模板（\ref{sec:vision-transforms-sharpening}节）。式\eqref{eq:vision-poisson-discrete}对$\Omega$中每个像素写一个方程，构成一个稀疏线性方程组。它的系数矩阵对称；若$\Omega$的每个连通部分都与边界相邻，矩阵还是正定的，因为对应的二次型等于$\Omega$内相邻像素差的平方和加上与边界相邻像素值的平方和，只有全零向量能使它为0。对称正定的稀疏方程组可以用共轭梯度法求解，可参考\BookChapterRef{chap:linear-algebra-and-matrix-analysis}{卷一“线性空间”章}；原论文用的是带逐次超松弛的Gauss–Seidel迭代或V循环多重网格。Gauss–Seidel迭代按扫描顺序逐个更新未知量，每次都使用当前最新的邻居值，见算法\ref{alg:vision-poisson-gauss-seidel}。

\begin{algorithm}[泊松图像编辑的Gauss–Seidel求解]
 \label{alg:vision-poisson-gauss-seidel}
 \begin{algorithmic}[1]
  \Require 目标图像$f^*$，源图像$g$，区域$\Omega$，松弛系数$\omega\in[1,2)$，扫描次数$T$
  \Ensure 区域内的结果$\{f_p\}_{p\in\Omega}$
  \State 初始化：对$p\in\Omega$令$f_p\gets g_p$；对相邻像素令$v_{p,q}\gets g_p-g_q$
  \For{$t=1,\ldots,T$}
   \For{按扫描顺序取$p\in\Omega$}
    \State $u\gets\Bigl(\sum_{q\in N_p\cap\Omega}f_q+\sum_{q\in N_p\cap\partial\Omega}f^*_q+\sum_{q\in N_p}v_{p,q}\Bigr)\big/|N_p|$
    \State $f_p\gets(1-\omega)f_p+\omega u$
   \EndFor
  \EndFor
  \State \Return $\{f_p\}_{p\in\Omega}$，区域外取$f^*$
 \end{algorithmic}
\end{algorithm}

用一维的例子可以手算全过程，见图\ref{fig:vision-poisson-1d}。像素编号0到5，$\Omega=\{1,2,3,4\}$，目标图像在两个边界像素上的值为$f^*_0=10$、$f^*_5=20$；源图像为$g=(0,0,4,8,4,0)$，像一个小鼓包。一维时每个像素只有左右两个邻居，式\eqref{eq:vision-poisson-discrete}给出
\begin{equation*}
 \begin{aligned}
 2f_1-f_2&=10+(0-0)+(0-4)=6,\\
 2f_2-f_1-f_3&=(4-0)+(4-8)=0,\\
 2f_3-f_2-f_4&=(8-4)+(8-4)=8,\\
 2f_4-f_3&=20+(4-8)+(4-0)=20\text{。}
 \end{aligned}
\end{equation*}
解为$f=(12,18,24,22)$，代回四个方程可以逐一验证。直接粘贴得到$(10,0,4,8,4,20)$，两端分别跳变$-10$和$+16$；泊松结果的相邻差为$(2,6,6,-2,-2)$，源图像的相邻差为$(0,4,4,-4,-4)$，每一处都多了2。边界上两幅图的总落差是$(20-10)-(0-0)=10$，被平均分摊到5个相邻差上，鼓包的形状完整保留，只是整体被抬高并倾斜。修正量$\tilde f=f-g=(12,14,16,18)$在10与20之间线性变化，这正是一维的$\Delta\tilde f=0$。若用$\omega=1$的Gauss–Seidel从$f=g$出发，第一遍扫描依次得到$f_1=(4+10-4)/2=5$、$f_2=(5+8+0)/2=6.5$、$f_3=(6.5+4+8)/2=9.25$、$f_4=(9.25+20+0)/2=14.625$，离解还很远：边界信息每遍只向内传播有限的距离，区域越大需要的扫描越多，这就是原论文同时采用多重网格的原因。

\begin{figure}[htbp]
 \centering
 % 一维泊松融合教学例：像素0-5，Omega={1,2,3,4}，边界值f*_0=10、f*_5=20，
% 源图像g=(0,0,4,8,4,0)；直接粘贴(10,0,4,8,4,20)，泊松解(10,12,18,24,22,20)。宽约93mm。
\begin{tikzpicture}[x=13mm,y=1.6mm,font=\figurefont\small]
 \fill[BookPaper] (.5,-1) rectangle (4.5,27);
 \node[anchor=north west,text=BookMuted] at (.55,27) {区域$\Omega$};
 \draw[BookMuted,line width=.7pt,->] (-.3,0)--(5.5,0) node[right] {像素};
 \draw[BookMuted,line width=.7pt,->] (-.3,0)--(-.3,28) node[above] {值};
 \foreach \y in {10,20} {
  \draw[BookRule,line width=.5pt] (-.3,\y)--(5.3,\y);
  \node[anchor=east] at (-.32,\y) {$\y$};
 }
 \node[anchor=east] at (-.32,0) {$0$};
 \foreach \x in {0,...,5} {
  \draw[BookMuted,line width=.5pt] (\x,0)--(\x,-.8);
  \node[below] at (\x,-.8) {$\x$};
 }
 \draw[BookBlue,dashed,line width=1pt] plot[mark=*,mark size=1.3pt,
  mark options={solid}] coordinates {(0,10)(1,0)(2,4)(3,8)(4,4)(5,20)};
 \draw[BookTeal,line width=1pt] plot[mark=*,mark size=1.3pt] coordinates {
  (0,10)(1,12)(2,18)(3,24)(4,22)(5,20)};
 \fill[BookMuted] (-.08,9.2) rectangle (.08,10.8);
 \fill[BookMuted] (4.92,19.2) rectangle (5.08,20.8);
 \draw[BookBlue,dashed,line width=1pt] (.2,-8.5)--(.6,-8.5);
 \node[anchor=west] at (.65,-8.5) {直接粘贴};
 \draw[BookTeal,line width=1pt] (2.2,-8.5)--(2.6,-8.5);
 \node[anchor=west] at (2.65,-8.5) {泊松融合};
 \fill[BookMuted] (4.05,-9.5) rectangle (4.25,-7.5);
 \node[anchor=west] at (4.3,-8.5) {边界值};
\end{tikzpicture}

 \caption{一维泊松融合的教学例。灰色方块是目标图像在边界像素0与5上的值，蓝色折线是直接粘贴源图像的结果，在两端各有一次跳变；青色折线是泊松融合的结果，两端与目标图像相接，内部保留了源图像的鼓包形状，每个相邻差都比源图像多2。数值为教学构造。}
 \label{fig:vision-poisson-1d}
\end{figure}

引导场不一定取源图像的梯度。把带孔洞或半透明的物体贴到有纹理的墙上时，可以在每对相邻像素上比较两幅图的变化，取较强的一个：$|f^*_p-f^*_q|>|g_p-g_q|$时取$v_{p,q}=f^*_p-f^*_q$，否则取$g_p-g_q$。原论文称之为混合梯度，它让墙面纹理透过字母之间的空隙保留下来。

泊松图像编辑成立的前提是源区域与目标区域在边界附近内容相容。若$\Omega$的边界横穿一条强边缘，例如桌沿，边界上的落差会沿着修正量扩散进花瓶，形成一片色斑；被贴物体的颜色总会被周围带偏，无法同时保持“边界无缝”和“颜色忠实于源图像”；它也不调整透视和光照方向，更不会为花瓶补上投在桌面的影子。“新内容在边界处与周围一致”这一要求在学习方法中仍然存在，只是不再显式解方程，而由修补模型读取可见区域、在训练中学会满足。

\subsubsection{Criminisi样例修补}
\label{sec:vision-generation-exemplar-inpainting}

泊松编辑需要一张现成的源图像。若请求只是“把杯子去掉”，杯子移除后留下的区域没有任何源图像可贴，只能用同一幅图中的桌面和背景补上。较早的修补方法沿着边界上的等照度线（亮度相等的曲线）把已知值扩散进来，适合细长的划痕，但填补大区域时会把纹理抹平；纹理合成能补出逼真的纹理，却不理会桌沿这样的直线结构\scite{vision-criminisi2004}。Criminisi等人的方法把纹理合成的采样方式与一个决定填充顺序的优先级结合起来，下面先看它所继承的纹理合成。

Efros与Leung把纹理看作马尔可夫随机场：给定一个像素周围$w\times w$窗口内的邻域，这个像素与图像其余部分无关\scite{vision-efros-leung1999}。窗口宽度$w$是唯一由用户设定的参数，应与纹理中最大的规则结构相当。合成一个像素$p$时，它的窗口$\omega(p)$中只有部分像素已知，算法在样例图中寻找最相似的窗口：距离$d$是两个窗口在已知像素上的平方差之和，用二维高斯核加权，使靠近中心的像素更重要，再除以已知像素的个数。设最小距离为$d_{\mathrm{best}}$，所有满足$d<(1+\epsilon)d_{\mathrm{best}}$的窗口组成候选集，原论文取$\epsilon=0.1$；从候选集中随机取一个窗口，用它的中心像素作为$p$的值。整幅纹理从一个随机取自样例的$3\times3$种子向外一层层生长，修补空洞时则从空洞边缘向内生长，见算法\ref{alg:vision-efros-leung}。

\begin{algorithm}[Efros–Leung非参数纹理合成]
 \label{alg:vision-efros-leung}
 \begin{algorithmic}[1]
  \Require 样例图$I_{\mathrm{smp}}$，待填像素集合，窗口宽度$w$，阈值系数$\epsilon=0.1$
  \Ensure 填好的图像
  \While{仍有待填像素}
   \State 取紧邻已知区域的一层待填像素
   \For{这一层中的每个像素$p$}
    \State 对样例图中每个窗口，按已知像素计算高斯加权、按已知像素数归一化的平方差$d$
    \State $d_{\mathrm{best}}\gets$最小的$d$；候选集$\gets\{$窗口$:d<(1+\epsilon)d_{\mathrm{best}}\}$
    \State 从候选集中随机取一个窗口，把它的中心像素值赋给$p$
   \EndFor
  \EndWhile
 \end{algorithmic}
\end{algorithm}

这一步为什么不会像回归那样取平均，可以用一维例子看清。设样例是一行像素$(0,1,1,0,1,1,0,1,2)$，窗口只看左边两个像素，要合成上下文为$(0,1)$之后的像素。样例中$(0,1)$出现在三处，后面分别跟着1、1、2，三处都完全匹配，$d_{\mathrm{best}}=0$，候选集就是这三处（$d_{\mathrm{best}}=0$时只保留完全匹配）。于是以$2/3$的概率取1、$1/3$的概率取2，每次取到的都是样例中真实出现过的值；若取三者的平均$4/3$，得到的是样例中从未出现的灰度。

逐像素合成很慢，且每个像素独立选择，大结构容易断裂。Efros与Freeman的图像缝合改为以$B\times B$的块为单位：相邻块之间重叠$B/6$的宽度，新块从样例中随机选取，要求它在重叠区的误差与最佳块相差不超过最佳误差的0.1倍；两块之间的接缝不取直线，而是沿重叠区误差最小的路径切开\scite{vision-quilting2001}。设重叠区第$i$行、第$j$列的平方差为$e_{i,j}$，累积误差按动态规划递推
\begin{equation*}
 E_{i,j}=e_{i,j}+\min\bigl(E_{i-1,j-1},E_{i-1,j},E_{i-1,j+1}\bigr),
\end{equation*}
最后一行$E$的最小值给出最小误差路径的终点，再逐行回溯。动态规划的一般原理可参考\BookChapterRef{chap:graph-optimization}{卷一“图”部分的“图上的约束优化”章}。以3行2列的重叠区为例，取$e=\left(\begin{smallmatrix}1&4\\3&1\\2&5\end{smallmatrix}\right)$，则第一行$E=(1,4)$，第二行$E=(3+1,1+1)=(4,2)$，第三行$E=(2+2,5+2)=(4,7)$；终点在第一列，回溯得到逐行的列号$(1,2,1)$，路径误差为$1+1+2=4$。

纹理合成把空洞当作一片纹理，按同心层的顺序填补，桌沿这样的直线结构在填到中间时往往错位。Criminisi等人指出，问题出在填充顺序：应当优先填补那些既有较多可靠邻居、又正好位于强边缘延长线上的位置。设待填区域为$\Omega$，已知的源区域为$\Phi$，二者之间的分界线称为填充前沿$\delta\Omega$。对前沿上的每个点$p$，以它为中心取一个块$\Psi_p$，原论文默认$9\times9$，实际应略大于纹理中最大的可辨元素。$p$的优先级定义为
\begin{equation}
 P(p)=C(p)D(p),\quad
 C(p)=\frac{\sum_{q\in\Psi_p\cap(I-\Omega)}C(q)}{|\Psi_p|},\quad
 D(p)=\frac{|\nabla I_p^{\perp}\cdot\bm{n}_p|}{\alpha}\text{。}
 \label{eq:vision-criminisi-priority}
\end{equation}
置信项$C(p)$是块内已知像素的置信度之和除以块的面积$|\Psi_p|$，初始化时$\Omega$内的像素置信度为0，其余为1；它偏向那些被已知像素包围较多的位置。数据项$D(p)$中，$\nabla I_p^{\perp}$是与梯度垂直、沿等照度线方向的向量，长度等于梯度幅值，$\bm{n}_p$是前沿在$p$处的单位法向，$\alpha$是归一化常数，灰度图取255；等照度线越强、越垂直地撞上前沿，$D(p)$越大，于是边缘被优先沿原方向延伸进空洞。选出优先级最高的$\hat p$以后，在源区域中找与$\Psi_{\hat p}$最相似的块$\Psi_{\hat q}$，距离是两块已知像素在CIE Lab颜色空间中的平方差之和；把$\Psi_{\hat q}$中对应位置的像素复制到$\Psi_{\hat p}$的未知像素上，并把这些新填像素的置信度设为$C(\hat p)$，见算法\ref{alg:vision-criminisi}。

\begin{algorithm}[Criminisi样例修补]
 \label{alg:vision-criminisi}
 \begin{algorithmic}[1]
  \Require 图像$I$，待填区域$\Omega$，块大小（默认$9\times9$）
  \Ensure 填好的图像
  \State 对$p\in\Omega$令$C(p)\gets0$，对其余像素令$C(p)\gets1$
  \While{$\Omega\neq\varnothing$}
   \State 找出填充前沿$\delta\Omega$，按式\eqref{eq:vision-criminisi-priority}计算前沿上每点的$P(p)$
   \State $\hat p\gets\arg\max_{p\in\delta\Omega}P(p)$
   \State 在源区域中找$\Psi_{\hat q}=\arg\min_{\Psi_q\subset\Phi}d(\Psi_{\hat p},\Psi_q)$，$d$只在$\Psi_{\hat p}$的已知像素上计算
   \State 把$\Psi_{\hat q}$的对应像素复制到$\Psi_{\hat p}\cap\Omega$
   \State 对新填像素$q$令$C(q)\gets C(\hat p)$，并把它们移出$\Omega$
  \EndWhile
 \end{algorithmic}
\end{algorithm}

用两个前沿点比较优先级。点$a$位于桌沿的延长线上，块内81个像素中已知54个，$C(a)=54/81\approx0.667$；桌沿的梯度幅值为120，等照度线垂直撞上前沿，$D(a)=120/255\approx0.471$，$P(a)\approx0.314$。点$b$位于平坦的桌面上，块内已知63个，$C(b)\approx0.778$；梯度幅值只有10，$D(b)\approx0.039$，$P(b)\approx0.030$。尽管$a$周围已知像素更少，它的优先级仍高出约10倍，桌沿会先被延伸进空洞。这些数值是教学构造。

这种贪心的逐块复制有明确的边界。它只能使用图像中已有的块，空洞里若应出现一个图中别处没有的物体，例如被杯子挡住的狗爪，就无法补出；某一步复制错误后，新填像素会成为后续匹配的依据，错误会沿前沿继续传播；块太小抓不住纹理元素，块太大又难以找到匹配。每次填块都要在整个源区域中搜索最相似的块，这一步的代价由下面的PatchMatch处理。

\subsubsection{PatchMatch}
\label{sec:vision-generation-patchmatch}

纹理合成、图像缝合和样例修补都反复执行同一个操作：对一个区域中的每个块，在另一个区域中找最相似的块。Barnes等人把这种对应关系称为\term{最近邻域场}（nearest-neighbor field, NNF）：对图像$A$中每个块的坐标$\bm{a}$，记它在图像$B$中最相似块的坐标为$\bm{b}$，则$f(\bm{a})=\bm{b}-\bm{a}$是一个偏移量，所有偏移量排成一张与$A$同尺寸的二维数组\scite{vision-patchmatch2009}。若区域含$M$个块位置、每块$m$个像素，穷举搜索的代价是$O(mM^2)$。百万像素的图像用$7\times7$的块，约为$49\times10^{12}$次像素运算，远不能交互使用；kd树等近似最近邻结构在高维的块空间中仍然慢，还要占用较多内存。

PatchMatch依据三点观察。其一，块空间的维数很高，但偏移量只是二维的，在偏移空间中搜索更省时省内存。其二，自然图像是连贯的：若某个块在$B$中的最佳匹配偏移了$\bm{v}$，它右边的块很可能也偏移$\bm{v}$，所以好的偏移可以传给邻居。其三，在大量随机猜测中，总有一部分恰好猜对。设$A$与$B$都含$M$个块位置，随机初始化时某个指定位置猜对的概率只有$1/M$，但至少有一个位置猜对的概率是$1-(1-1/M)^M$，$M$很大时趋于$1-1/e\approx0.63$；若把落在正确偏移附近$C$个像素内的猜测也算作好的猜测，这个概率是$1-(1-C/M)^M\approx1-e^{-C}$，$C=5$时约为0.993。只要连贯区域中有一个位置猜对，传播就能把它扩散到整个区域。

算法先给每个块随机分配一个偏移，然后反复迭代，见算法\ref{alg:vision-patchmatch}。每次迭代按扫描顺序（从左到右、从上到下）访问每个块，对它依次做传播和随机搜索。传播时，用$D(\bm{v})$表示$A$中$(x,y)$处的块与$B$中$(x,y)+\bm{v}$处的块之间的距离，在$\{f(x,y),f(x-1,y),f(x,y-1)\}$中取使$D$最小的偏移；偶数次迭代按逆扫描顺序进行，改用$f(x+1,y)$与$f(x,y+1)$，使好的偏移也能向左、向上传播。随机搜索在当前偏移$\bm{v}_0=f(x,y)$周围按指数缩小的范围试探
\begin{equation}
 \bm{u}_i=\bm{v}_0+w\alpha^i\bm{R}_i,\qquad i=0,1,2,\ldots
 \label{eq:vision-patchmatch-search}
\end{equation}
其中$\bm{R}_i$在$[-1,1]\times[-1,1]$中均匀随机，$w$是最大搜索半径，原论文取图像的最大边长，$\alpha=1/2$；半径$w\alpha^i$小于1个像素时停止，候选落到$B$外时截到边界上，凡是更好的候选都替换当前偏移。原论文的结果都只迭代4到5次，此时偏移场几乎不再变化。

\begin{algorithm}[PatchMatch近似最近邻域场]
 \label{alg:vision-patchmatch}
 \begin{algorithmic}[1]
  \Require 图像$A,B$，块距离$D$，迭代次数$T$（4到5），搜索半径$w$，缩小比例$\alpha=1/2$
  \Ensure 偏移场$f$
  \State 对$A$中每个块位置，从$B$的范围内均匀随机抽取$f(x,y)$
  \For{$t=1,\ldots,T$}
   \If{$t$为奇数}
    \State 扫描顺序为从左到右、从上到下，邻居集合$S(x,y)=\{(x-1,y),(x,y-1)\}$
   \Else
    \State 扫描顺序反向，$S(x,y)=\{(x+1,y),(x,y+1)\}$
   \EndIf
   \For{按扫描顺序取每个块位置$(x,y)$}
    \State 传播：$f(x,y)\gets\arg\min_{\bm{v}\in\{f(x,y)\}\cup\{f(\bm{s}):\bm{s}\in S(x,y)\}}D(\bm{v})$
    \State 随机搜索：对$i=0,1,\ldots$直到$w\alpha^i<1$，按式\eqref{eq:vision-patchmatch-search}生成$\bm{u}_i$，若$D(\bm{u}_i)<D(f(x,y))$则$f(x,y)\gets\bm{u}_i$
   \EndFor
  \EndFor
  \State \Return $f$
 \end{algorithmic}
\end{algorithm}

用一维例子手算传播。设$A$中位置0到3的四个块，在$B$中的真实匹配都在10到13，即正确偏移都是$+10$；随机初始化得到$(+3,+10,-1,+5)$，只有位置1碰巧猜对。第一次迭代从左到右：位置0没有左邻，假设随机搜索没有命中；位置1已正确；位置2比较自身的$-1$与左邻的$+10$，后者距离为0，取$+10$；位置3同理取$+10$。第二次迭代从右到左：位置3、2、1不变，位置0比较自身的$+3$与右邻的$+10$，取$+10$。两次迭代后全部正确，而初始时只有一个位置猜对。随机搜索的候选数也可以算出：若$w=1000$，半径依次为$1000,500,\ldots,1000/2^9\approx1.95$，下一个$1000/2^{10}<1$，所以每个块每次迭代试探10个候选。图\ref{fig:vision-patchmatch-steps}画出了初始化和两次传播后的偏移场。

\begin{figure}[htbp]
 \centering
 % PatchMatch一维教学例：A中位置0-3的正确偏移均为+10（匹配B中10-13）。
% 初始化(+3,+10,-1,+5)；第一次正向传播后(+3,+10,+10,+10)；第二次逆向传播后全部为+10。
% 三个面板上下排列，每格5mm，宽约107mm、高约108mm。
\begin{tikzpicture}[x=5mm,y=1mm,font=\figurefont\small,
 cellA/.style={draw=BookBlue,fill=BookBlue!10,line width=.6pt},
 cellB/.style={draw=BookMuted,fill=white,line width=.5pt},
 good/.style={-{Stealth[length=1.6mm]},BookTeal,line width=.9pt},
 bad/.style={-{Stealth[length=1.6mm]},BookBlue,dashed,line width=.8pt},
 prop/.style={-{Stealth[length=1.4mm]},BookAmber,line width=.8pt}]
 \foreach \panel/\ytop/\title/\offs in {
   0/0/{初始化}/{3,10,-1,5},
   1/-34/{第一次迭代：正向传播}/{3,10,10,10},
   2/-68/{第二次迭代：逆向传播}/{10,10,10,10}} {
  \node[anchor=west] at (-4.6,\ytop+5) {\title};
  \node[anchor=east] at (-.2,\ytop-3) {$A$};
  \node[anchor=east] at (-.2,\ytop-19) {$B$};
  \foreach \i in {0,...,3} {
   \draw[cellA] (\i,\ytop-6) rectangle (\i+1,\ytop);
   \node at (\i+.5,\ytop-3) {\footnotesize\i};
  }
  \foreach \j in {0,...,15} {
   \draw[cellB] (\j,\ytop-22) rectangle (\j+1,\ytop-16);
   \node[text=BookMuted] at (\j+.5,\ytop-19) {\scriptsize\j};
  }
  \foreach \o [count=\k from 0] in \offs {
   \pgfmathtruncatemacro{\tgt}{\k+\o}
   \ifnum\o=10
    \draw[good] (\k+.5,\ytop-6.3) -- (\tgt+.5,\ytop-15.7);
   \else
    \draw[bad] (\k+.5,\ytop-6.3) -- (\tgt+.5,\ytop-15.7);
   \fi
  }
 }
 \draw[prop] (1.6,-34+0.8) to[bend left=50] (2.4,-34+0.8);
 \draw[prop] (2.6,-34+0.8) to[bend left=50] (3.4,-34+0.8);
 \draw[prop] (1.4,-68+0.8) to[bend right=50] (0.6,-68+0.8);
 \draw[good] (3,-96)--(4,-96); \node[anchor=west] at (4.1,-96) {正确偏移};
 \draw[bad] (7.5,-96)--(8.5,-96); \node[anchor=west] at (8.6,-96) {错误偏移};
 \draw[prop] (12,-96) to[bend left=40] (13,-96); \node[anchor=west] at (13.1,-96) {传播方向};
\end{tikzpicture}

 \caption{PatchMatch传播过程的一维教学例。每个面板上排是$A$中位置0至3的块，下排是$B$中的位置，箭头从$A$的块指向当前偏移所指的$B$中位置，实线为正确偏移$+10$，虚线为错误偏移；弯箭头表示本次迭代中好偏移的传播方向。初始化时只有位置1碰巧正确；第一次迭代按扫描顺序把$+10$传给位置2和3，第二次迭代逆向扫描再传给位置0。图中省略了随机搜索。数值为教学构造。}
 \label{fig:vision-patchmatch-steps}
\end{figure}

在原论文的修补中，空洞记为$T$，图像其余部分记为$S$，修补按类似期望最大化的方式交替进行：先用PatchMatch为$T$中每个块在$S$中找最近邻，每个匹配块把自己的像素“投给”它覆盖的位置；再让每个像素取所得投票的平均值，得到新的$T$。为避免陷入差的局部解，计算从低分辨率开始，逐级放大并在每级继续迭代。论文在最大约200万像素的测试图像上验证了收敛性，并报告CPU实现比带PCA降维的kd树快20到100倍，典型参数下的时间复杂度为$O(mM\log M)$，额外内存为$O(M)$，可以交互使用。

PatchMatch只加速了搜索，没有改变样例方法的根本局限：新内容仍只能来自已有的块。它的偏移场也只按块距离优化，不包含“这是一只杯子”之类的语义，空洞里应出现的内容若在图中找不到相似的块，结果只能是貌似合理的纹理拼贴。学习方法中的注意力也计算一个位置与所有位置的相似度，再按相似度从其他位置取值，可以看成一种可学习的软性最近邻；而要生成原图中不存在的花瓶，则需要下面以文字为条件的生成模型。

\subsubsection{GLIDE修补模型}
\label{sec:vision-generation-glide-inpainting}

最简单的修补做法是：用普通文字生成模型采样，在每一步把掩码外的像素替换为原图相应噪声水平的像素。这样能强制保留已知区域，但模型在训练中从未见过“部分像素已知”的输入，它不知道已知区域里有什么，生成的花瓶常与周围光照和边界不协调。GLIDE修补模型的动机是让模型在每一步都能直接看到清晰的保留区域，从而学会利用上下文\scite{vision-glide2022}。

主要想法是把可见图像和掩码作为额外输入通道，对\ref{sec:vision-models-glide}节的文字条件GLIDE作修补专用的微调。模型在每个去噪步同时读到当前带噪图像、清晰的保留区域和掩码，因此可以根据桌面纹理、光照方向来生成缺失区域。

架构上，GLIDE基础模型在$64\times64$的像素空间工作。修补模型的输入是三部分沿通道拼接：当前带噪图像（3通道）、可见图像$\bm{V}=\bm{M}\odot\bm{X}$（3通道）和掩码$\bm{M}$（1通道），共$64\times64\times7$。只有第一层卷积的输入通道从3增加到7，新增4个通道对应的权重初始化为零，其余权重从文字生成检查点继承，所以微调开始时模型的行为与原模型相同。文字仍经Transformer编码后接入去噪网络。上采样模型把分辨率提到$256\times256$，它读取完整的低分辨率图像，以及高分辨率下的未遮挡区域。掩码通道不可省略：可见图像中被遮住的位置填0，而真实图像中也可能有黑色像素，只有掩码能区分二者。

各组件的规模见表\ref{tab:vision-glide-inpaint-config}。基础模型沿用Dhariwal与Nichol的ImageNet $64\times64$扩散模型（ADM）的结构\scite{vision-adm2021}，只把宽度扩大到512通道，视觉部分约23亿参数；文字先编码为$K$个词元，送入一个24个残差块、宽2048的Transformer，约12亿参数。这个Transformer的输出有两种用法：最后一个词元的向量代替ADM中的类别嵌入，最后一层的$K$个词元向量则分别投影到ADM每个注意力层的维数，拼接到该层的注意力上下文中，让图像特征可以注意到文字。GLIDE论文没有逐层列出ADM的结构，各分辨率的残差块数与注意力位置以ADM论文为准，本章不再展开。

\begin{table}[htbp]
 \centering
 \small
 \begin{tabularx}{\linewidth}{@{}p{24mm}Xp{20mm}@{}}
  \toprule
  组件 & 结构与输入输出 & 规模\\
  \midrule
  文字编码器 & Transformer，24个残差块，宽2048；输出$K$个词元向量 & 约12亿参数\\
  基础扩散模型 & ADM结构，宽512通道；$64\times64$像素空间 & 视觉部分约23亿参数\\
  修补模型的输入层 & 第一层卷积由3通道扩为7通道，新增4通道的权重初始化为0 & —\\
  上采样扩散模型 & $64\to256$，ADM上采样结构，基础通道384；文字编码器宽1024 & 约15亿参数\\
  训练 & 基础模型250万次迭代、批大小2048；上采样模型160万次迭代、批大小512；再以20\%空文字微调以支持无分类器引导 & —\\
  修补微调 & 随机擦除区域，以可见图像与掩码为附加输入；完整模型的修补微调步数论文未给出 & —\\
  采样 & 修补样例用100个扩散步，其余展示样例用150步 & —\\
  \bottomrule
 \end{tabularx}
 \caption{GLIDE及其修补模型的配置，数值取自GLIDE论文。基础模型与文字编码器合计约35亿参数；论文公开的小规模过滤版本在110万次预训练之后，另用50万次迭代同时微调无分类器引导与修补。}
 \label{tab:vision-glide-inpaint-config}
\end{table}

训练数据就是原有的图文对，无需额外标注。每个样本随机擦去图像的一些区域得到$\bm{M}$，构造$\bm{V}$，再对完整图像加噪，以完整图像的噪声预测误差微调模型。由于每次抽到的掩码不同，模型学会在各种可见范围下补全。推理时固定用户给出的掩码、原图和文字“花瓶”，从噪声出发反复去噪；可见条件在整个采样中保持不变，只有带噪图像随步更新。

这里要分清训练目标和使用目标。训练时，模型补全的是被擦去的真实内容；使用时，文字要求的花瓶原图中并不存在，这属于创作，模型借助的是训练中学到的“依上下文补全”的能力。若只是要恢复被遮挡的杯子，那是\ref{sec:vision-restoration-inpainting}节的复原目标。掩码条件不保证掩码外的每个像素都不变，因为生成和上采样都在全图上进行。需要严格保持时，可以按掩码把原图回填，这是一个额外的后处理步骤，回填后要重新检查边界。

\subsubsection{InstructPix2Pix}
\label{sec:vision-generation-instruct-pix2pix}

用户常常只给一句“把杯子换成花瓶”，不画掩码，也不写完整的新场景描述。这时需要一个读入原图和编辑指令、直接输出编辑结果的模型。训练这样的模型需要大量“原图、指令、编辑后图像”三元组，而人工制作这样的配对几乎不可行。InstructPix2Pix的动机就是自动合成这些训练数据，再用它们训练一个指令编辑模型\scite{vision-instruct-pix2pix2023}。

主要想法分两步。第一步合成数据：用语言模型从图像描述生成编辑指令和编辑后的描述，再用文字生成模型把编辑前后两段描述分别生成图像。第二步训练一个以原图和指令为条件的扩散模型。合成数据的关键困难是，独立生成的两幅图会在姿态、背景上全面不同，不是一对编辑前后的图像。InstructPix2Pix用Prompt-to-Prompt方法解决：两幅图使用相同的初始噪声和相同的采样噪声，并在前一部分去噪步中让第二幅图沿用第一幅图的注意力权重，使两幅图的布局保持一致，只在描述不同的地方变化\scite{vision-prompt2prompt2022}。Prompt-to-Prompt原本按编辑类型替换交叉注意力；InstructPix2Pix的附录说明，它改为对所有编辑统一替换第二幅图的自注意力权重。

数据流程的细节如下。先人工撰写700条三元组（输入描述、编辑指令、输出描述），用默认训练参数把GPT-3微调一轮；再让微调后的GPT-3为大量图像描述生成指令和新描述。对每对描述，用Stable Diffusion配合Prompt-to-Prompt生成100对图像，共享注意力的步数比例在0.1到0.9之间随机抽取；然后用CLIP相似度筛选，要求两图的CLIP图像相似度超过0.75、图文相似度超过0.2，以及“图像变化方向与描述变化方向”的方向相似度超过0.2，通过的配对按方向相似度排序，每对描述最多保留4对。生成图像用Stable Diffusion v1.5的指数滑动平均权重、Euler祖先采样器和100个去噪步。最终得到超过45万个训练样本。共享注意力减少了无关变化，筛选去掉了明显失败的配对，但并不保证每条指令都被正确执行。

模型以Stable Diffusion v1.5初始化。原图经冻结的自编码器编码为潜表示$\bm{z}_I$；训练时目标图像编码后加噪得到$\bm{z}^{(t)}$；两者都是$h'\times w'\times4$，沿通道拼接为8通道输入去噪U-Net的第一层，新增4个通道的权重初始化为零。指令$c_T$经文字编码器和交叉注意力进入网络，不拼进像素通道。训练目标为
\begin{equation}
 L=\mathbb{E}\!\left[
 \left\lVert\bm{\epsilon}-
 \epsilon_{\bm{\theta}}
   \bigl(\bm{z}^{(t)},t,\bm{z}_I,c_T\bigr)
 \right\rVert_2^2\right]\text{，}
 \label{eq:vision-instruct-training}
\end{equation}
其中$\bm{\epsilon}$是加到目标潜表示上的噪声。原图条件在整个训练和使用中都保持干净，只有目标在加噪。为支持下面的双条件引导，训练时以5\%的概率只丢弃图像条件、5\%只丢弃文字条件、5\%同时丢弃二者。训练配置见表\ref{tab:vision-instruct-config}。论文说明全部可用权重都从预训练检查点初始化，其余训练设置沿用公开的Stable Diffusion代码库，但没有单独说明自编码器和文字编码器是否参与更新。

\begin{table}[htbp]
 \centering
 \small
 \begin{tabularx}{\linewidth}{@{}p{24mm}Xp{20mm}@{}}
  \toprule
  组件 & 结构与输入输出 & 规模\\
  \midrule
  去噪U-Net & Stable Diffusion v1.5的U-Net（逐层结构见表\ref{tab:vision-controlnet-structure}），从其指数滑动平均权重初始化 & 论文未给出\\
  输入层 & 第一层卷积读取带噪目标潜表示与原图潜表示，共8通道；新增4通道的权重初始化为0 & —\\
  条件接入 & 指令经原有的文字条件机制（交叉注意力）进入；原图只经输入通道拼接进入 & —\\
  训练数据 & 超过45万个由GPT-3与Stable Diffusion合成的“原图、指令、编辑后图像”样本 & —\\
  训练 & $256\times256$，1万步，总批大小1024，学习率$10^{-4}$且无预热；随机水平翻转，随机缩放到256至288再裁回256；8张40GB A100约25.5小时 & —\\
  推理 & $512\times512$，Euler祖先采样器100步，每步三次前向；单张A100约9秒一幅 & —\\
  \bottomrule
 \end{tabularx}
 \caption{InstructPix2Pix的配置，数值取自原论文正文与附录。模型在$256\times256$上训练，论文报告它在推理时能推广到$512\times512$。}
 \label{tab:vision-instruct-config}
\end{table}

推理时输入原图和指令，从随机噪声开始采样，不需要目标图像。每步计算三种噪声预测：$\bm{\epsilon}_0=\epsilon_{\bm{\theta}}(\bm{z}^{(t)},t,\varnothing,\varnothing)$、$\bm{\epsilon}_I=\epsilon_{\bm{\theta}}(\bm{z}^{(t)},t,\bm{z}_I,\varnothing)$和$\bm{\epsilon}_{I,T}=\epsilon_{\bm{\theta}}(\bm{z}^{(t)},t,\bm{z}_I,c_T)$，再组合为
\begin{equation}
 \hat{\bm{\epsilon}}
 =\bm{\epsilon}_0
  +s_I(\bm{\epsilon}_I-\bm{\epsilon}_0)
  +s_T(\bm{\epsilon}_{I,T}-\bm{\epsilon}_I)\text{。}
 \label{eq:vision-instruct-guidance}
\end{equation}
这是无分类器引导在两个条件上的推广：第二项强调原图，第三项在已有原图条件下强调指令带来的变化。以两个分量示意，令$\bm{\epsilon}_0=(0.1,0.2)$、$\bm{\epsilon}_I=(0.2,0.1)$、$\bm{\epsilon}_{I,T}=(0.4,-0.1)$，取$s_I=1.5,s_T=5$，得$(1.25,-0.95)$；两项强度都取1时恢复$\bm{\epsilon}_{I,T}$。原论文指出$s_T$在5到10、$s_I$在1到1.5之间通常效果较好，并需按例调整：$s_I$越大越接近原图，$s_T$越大修改越强。每步三次前向使推理成本约为普通采样的三倍。

由于没有显式掩码，未提及的区域也可能改变，例如狗的颜色随杯子一起变化。数量、空间关系类的指令以及需要大幅改变几何的指令容易失败，这部分错误继承自合成数据：生成数据的文字生成模型和筛选器本身就不擅长这些情形。

五种单轮编辑方法的比较见表\ref{tab:vision-generation-editing-comparison}。固定任务为“把狗旁边的杯子换成花瓶”，评价口径为替换是否完成、杯子以外是否保持，以及新内容与周围是否协调。三种经典方法都不训练模型，新内容来自现成图像：泊松图像编辑需要一张花瓶照片，在保留区严格不变的前提下用梯度域求解消除接缝，代价是花瓶颜色被带偏；Criminisi样例修补与PatchMatch只能完成“去掉杯子”这一半，用图中已有的桌面补洞，前者规定填充顺序，后者把块搜索加速到可交互，二者都补不出图中没有的花瓶。GLIDE修补模型把新内容的来源换成以文字为条件的生成模型，掩码仍由用户给出，可见区域作为输入让模型学会与周围协调，代价是需要大规模图文数据训练，且掩码外不保证逐像素不变。InstructPix2Pix相对GLIDE修补模型又改变了两处设计：范围由用户掩码改为模型从指令中推断，条件接入由“可见图像加掩码”改为“完整原图潜表示”。它换来了只需一句话的使用方式，代价是失去对保留区域的显式约束，并依赖合成数据的质量。

\begin{table*}[tbp]
 \centering
 \small
 \begin{tabularx}{\linewidth}{@{}p{26mm}p{30mm}XXX@{}}
  \toprule
  方法 & 用户输入 & 新内容的来源 & 范围外的保持 & 典型失败\\
  \midrule
  泊松图像编辑 & 原图、区域、源图像 & 源图像的梯度，经求解泊松方程融合 & 区域外严格不变 & 边界横穿强边缘时出现色斑，被贴物体颜色偏移，不处理透视与阴影\\
  Criminisi样例修补 & 原图、待填区域 & 同一图中别处的图像块，按优先级逐块复制 & 区域外严格不变 & 图中没有的内容补不出，复制错误沿前沿传播\\
  PatchMatch & 原图、待填区域 & 同一图中的图像块，由随机化最近邻搜索与投票合成 & 区域外严格不变 & 与样例修补相同；只加速搜索，不含语义\\
  GLIDE修补模型 & 原图、掩码、对新内容的描述 & 文字条件扩散模型，读取可见区域与掩码 & 由掩码约束，可回填强制保持 & 掩码过小留下旧阴影、边界不协调\\
  InstructPix2Pix & 原图、编辑指令 & 以原图潜表示和指令为条件的扩散模型 & 无显式约束，靠$s_I$调节 & 未提及区域改变，数量与几何指令失败\\
  \bottomrule
 \end{tabularx}
 \caption{五种单轮编辑方法在同一替换任务下的比较。经典方法不需要训练，每次编辑的计算量随区域大小增长；两种学习方法需要预先在大规模数据上训练，单次编辑的代价是一次完整的扩散采样，InstructPix2Pix每步还需三次前向。}
 \label{tab:vision-generation-editing-comparison}
\end{table*}

\subsection{多轮编辑的状态维护}
\label{sec:vision-generation-kontext-history}

一次替换之后，用户还可能继续要求“把花瓶改成蓝色”“再把它移到狗的右边”。每一轮都有新的变化，同时已经接受的修改和始终不应变化的主体都要保持。若只把当前图像送入下一轮，模型能看到上一轮的结果，却无法恢复已被前几轮采样削弱的狗的斑纹；若每轮都从最初的原图开始，又会丢掉已接受的花瓶。因此多轮编辑需要维护一份状态：原始依据、每轮的已接受版本以及修改记录，并写明每一项是保持还是被本轮覆盖。这份状态由模型的输入能力和外部的版本管理共同完成，二者应分开描述。

模型一侧，需要一个能把已有图像作为上下文、按指令生成新图的接口，而且最好能在多轮中保持主体。FLUX.1 Kontext是这样的模型\scite{vision-kontext-2025}。它的动机是：InstructPix2Pix沿通道拼接原图和目标，要求二者尺寸相同，并需要为编辑单独训练；而Kontext希望用同一个模型完成文字生成和图像编辑，让条件图与目标图的尺寸可以不同。主要想法是沿序列方向拼接：上下文图像经冻结自编码器编码并分块后成为$n_y$个词元，待生成图像成为$n_x$个词元，二者拼成长度$n_x+n_y$的序列送入以流匹配训练的Transformer，文字词元经另一条流与之交换注意力。为区分两幅图中相同空间坐标的词元，三维旋转位置编码给目标词元赋$(0,h,w)$，给上下文词元赋$(1,h,w)$，首坐标只是图像来源编号。

训练从FLUX.1文字生成检查点开始，样本为“上下文图像$\bm{Y}$、指令$c$、目标图像$\bm{X}$”，纯文字生成样本令$\bm{Y}$为空。设$\bm{x}=E(\bm{X})$为目标潜表示，抽标准高斯噪声$\bm{\epsilon}$与$t\in(0,1)$，目标为
\begin{equation}
 \begin{aligned}
 \bm{z}(t)&=(1-t)\bm{x}+t\bm{\epsilon},\\
 L_{\mathrm{ctx}}&=
 \mathbb{E}\!\left[
 \left\lVert v_{\bm{\theta}}
  (\bm{z}(t),t,E(\bm{Y}),c)
   -(\bm{\epsilon}-\bm{x})\right\rVert_2^2
 \right]\text{。}
 \end{aligned}
 \label{eq:vision-kontext-flow}
\end{equation}
误差只计在目标词元上，上下文词元保持干净、不作为输出。这里$t=0$是干净目标、$t=1$是噪声，与\ref{sec:vision-models-flow-matching}节的方向相反，所以采样沿$t$递减积分。自编码器冻结，更新的是生成主干。论文没有公开全部关系数据构造与训练预算；公开的dev版本还经过引导蒸馏，推理步数与完整模型不同。

Kontext的上下文保持干净，并不意味着输出在对应区域逐像素不变：编码、全图生成和解码都会改变细节，模型也不会因为序列拼接而自动保存编辑历史。论文的形式可以扩展到多幅上下文图像，但所发布的版本主要以单幅上下文训练。所以外部版本管理不可省略：保存原图$\bm{X}^{(0)}$，接受杯子替换得到$\bm{X}^{(1)}$，花瓶变蓝得到$\bm{X}^{(2)}$，移动花瓶得到$\bm{X}^{(3)}$。每轮以当前已接受的图像作为单幅上下文，同时记录模型、提示、种子、尺寸和候选数。第三轮要同时检查花瓶的位置、是否仍为蓝色、狗的耳朵和斑纹、未要求改变的背景。若第三轮颜色退回，应回到$\bm{X}^{(2)}$重做，而不是把新结果当成下一轮的正确起点。

历史要求也有不同类型。ImgEdit把多轮编辑的评价分成三类：早期设定的全局要求在后续轮次中仍然有效，例如“所有结果保持木质纹理”；后一轮请求依赖前文指代，例如“把它改成蓝色”中的“它”；以及回到某个历史版本的操作\scite{vision-imgedit-2025}。这三类要求都依赖保存的状态，而不是单轮的“原图—结果”配对。随着轮数增加，编码、采样和解码误差会累积，即使某轮没有要求改动主体，主体也可能漂移；Kontext论文同样展示了多轮后出现伪影和指令遗漏的失败例。外部回退可以阻止错误继续传播，却补不回原始参考从未提供的细节。

\subsection{编辑结果的核验}
\label{sec:vision-generation-editing-check}

核验时应把原图、编辑结果、要求记录和允许变化的范围放在一起。杯子仍在，是目标未修改；花瓶出现在另一侧，是改错位置；狗的斑纹变化，是保留失败；花瓶与桌面之间出现断边，是边界问题。把这些错误合成一个“编辑质量”分数，就无法判断应该修改流程中的哪一步。

对于明确的保留掩码$\bm{M}$，用式\eqref{eq:vision-edit-preservation}计算保留区的平均变化。若直接返回原图，$E_{\mathrm{keep}}=0$，但杯子没有变成花瓶，所以这一项必须与修改完成度一起报告。若请求允许全图色调变化，逐像素的保留误差就要换成与要求相符的判断。泊松图像编辑是一个需要特别说明的例子：它在保留区内严格不变，$E_{\mathrm{keep}}=0$，但粘贴进来的花瓶整体颜色会被调整，这一变化要按“新内容是否保持源物体外观”另行检查。

为强制保持范围外的内容，可以合成$\bm{X}_{\mathrm{out}}=\bm{M}\odot\bm{X}+(1-\bm{M})\odot\bm{X}'$。它保证保留像素严格相等，却可能截断新花瓶的阴影；若使用羽化掩码，就要说明过渡带也会改变。多轮编辑时，每轮都应与最初的原图和保留要求对照，而不只与上一轮比较，否则身份漂移会逐轮积累而不被发现。缩小范围、调整条件强度和重新生成都要记录轮数与候选预算，最终成功的结果应连同得到它的成本一起报告。

\section{观测场景的新视角合成}
\label{sec:vision-generation-novel-view}

回到开篇的第四个请求：根据绕玩具狗拍摄的一圈照片，给出从正后方看的样子。\term{新视角合成}（novel view synthesis）根据已拍摄的视图，渲染同一场景在未拍摄相机位置下的图像。它与前面三类任务的根本区别在于，结果应由真实观测支持，而不是由生成先验补充。输入是多幅照片及其相机参数，输出是任意给定相机下的图像，难点在于从有限的二维观测中拟合一个能够解释所有视图的三维表示。本节按“组织相机与观测—拟合场景表示—渲染新视角并核验”的步骤展开。只有一幅图、需要靠先验补出背面的情形，属于\ref{chap:vision-multimodal-generation}章的多视角生成。

\subsection{相机与观测的组织}
\label{sec:vision-generation-cameras}

拟合场景之前，每幅照片都要有已知的相机参数。这里用到\ref{chap:vision-image-formation}章建立的针孔相机模型：内参矩阵$\bm{K}$（焦距、主点）决定像素与相机坐标系中方向的关系，外参即旋转$\bm{R}$与平移$\bm{t}$决定相机在世界中的位置和朝向，世界点$\bm{x}$投影到像素的齐次坐标为$\bm{K}(\bm{R}\bm{x}+\bm{t})$，再除以第三个分量（\ref{sec:vision-formation-camera-intrinsics}节）。把这个过程反过来，像素$(u,v)$对应一条从相机中心$\bm{o}=-\bm{R}^{\top}\bm{t}$出发、方向为$\bm{d}\propto\bm{R}^{\top}\bm{K}^{-1}(u,v,1)^{\top}$的射线，这是后面所有渲染计算的出发点。例如$\bm{K}$的焦距为500像素、主点为$(320,240)$，相机位于世界原点且不旋转，则像素$(570,240)$对应方向$(0.5,0,1)$，即向右偏离光轴约26.6度。

相机参数通常由运动恢复结构从照片本身估计（\ref{sec:vision-formation-multi-view}节）。它先在照片之间匹配局部特征，由两视图的对极几何得到相对位姿（\ref{sec:vision-formation-two-view}节），再逐步加入更多照片，用三角化得到三维点，最后用光束法平差同时调整所有相机和点，使三维点在各照片上的重投影误差最小。输出是每幅照片的$\bm{K},\bm{R},\bm{t}$和一组稀疏的三维点。NeRF对真实场景使用COLMAP运动恢复结构软件估计相机\scite{vision-nerf2020}，3D Gaussian Splatting还用它给出的稀疏点初始化场景。相机也可以用\ref{chap:vision-image-recognition}章“相机与场景几何的估计”中的前馈几何模型估计。多视图几何本身只给出相机和稀疏点；要从任意新相机渲染出逼真图像，还需要一个能表示所有可见表面外观的场景表示，这是下一小节的问题。

拍摄的覆盖范围决定哪些表面有多视角证据。绕玩具狗一圈、在两个高度上各拍若干张，可以覆盖大部分表面；若只拍了正面，背面在任何拟合中都没有观测支持。拍摄前还应留出一些相机作为测试视图，它们不参与拟合，用于后面的核验。原始NeRF和3D Gaussian Splatting都要求场景静止、曝光一致、相机准确：玩具被碰动、光照变化或标定误差，都会让同一空间点需要解释互相矛盾的像素，拟合过程可能用错误的密度或外观吸收这些矛盾，训练视图看起来变好，几何却更不准确。

\subsection{场景表示的拟合}
\label{sec:vision-generation-scene-fitting}

有了带相机的照片，需要决定用什么表示场景，才能从新相机渲染出图像。基于图像的渲染给出了最早的一种回答：不重建几何，直接把照片看成对光线的采样。在没有遮挡物的自由空间中，沿一条射线传播的光强不变，所以可以用射线与两个平行平面的交点$(u,v)$和$(s,t)$来索引它，全部射线构成一个四维函数，称为光场；新视角的每个像素对应一条射线，只需在已拍到的射线中查找并插值\scite{vision-lightfield1996}。这种做法不需要估计任何深度，但只有在相机排布得足够密集、相邻射线足够接近时，插值才可靠；只拍了一圈几十张照片时，大多数新射线离已拍射线很远，插值结果会出现重影。另一条经典路线先用多视图立体重建出网格，再把照片作为纹理贴上去，对表面清晰的物体很有效，但毛绒、半透明和细小结构难以得到准确的网格。

学习路线把场景表示成一组可优化的参数，要求渲染过程可微，从而用“渲染图与照片的差异”直接优化表示，不必先得到准确的表面。两种代表性的路线对“场景用什么表示、怎样渲染”给出了不同回答。NeRF用一个神经网络隐式表示连续的密度和颜色场，渲染时沿每条像素射线采样并积分；3D Gaussian Splatting用数百万个显式的三维高斯椭球表示场景，渲染时把它们投影到图像平面并按深度合成。两者都对每个场景单独优化，换一个场景要重新拟合；它们的差别主要体现在渲染速度、存储和细节表达上。

\subsubsection{NeRF}
\label{sec:vision-generation-nerf}

在NeRF之前，新视角合成多依赖显式的网格、点云或离散体素。网格需要先重建出准确的表面，对毛绒、半透明和细小结构很困难；离散体素的存储随分辨率的立方增长，细节受限。\term{神经辐射场}（neural radiance field, NeRF）的动机是用一个连续函数表示场景，让存储与分辨率无关，并用可微的体渲染从照片直接优化它\scite{vision-nerf2020}。

主要想法是：用一个多层感知机把三维位置$\bm{x}\in\mathbb{R}^3$和观察方向$\bm{d}$映射为体密度$\sigma(\bm{x})\geq0$和颜色$\bm{c}(\bm{x},\bm{d})\in[0,1]^3$。密度描述该位置有多少“物质”会截获光线，只依赖位置；颜色还依赖方向，以表达高光随视角的变化。渲染一个像素时，沿它的射线累加各处的颜色贡献，这种按密度与遮挡加权累加的计算称为\term{体渲染}（volume rendering）。整个渲染过程对网络参数可微，因此可以用照片像素直接监督。

先写出连续形式。相机像素对应射线$\bm{r}(s)=\bm{o}+s\bm{d}$，$\bm{o}$是相机中心，$s$是沿射线的距离。在近远界$s_{\mathrm{n}},s_{\mathrm{f}}$之间，像素颜色为
\begin{equation}
 \begin{aligned}
 \bm{C}(\bm{r})
 &=\int_{s_{\mathrm{n}}}^{s_{\mathrm{f}}}
    T(s)\sigma(\bm{r}(s))
    \bm{c}(\bm{r}(s),\bm{d})\,\mathrm{d}s,\\
 T(s)&=\exp\!\left(
     -\int_{s_{\mathrm{n}}}^{s}\sigma(\bm{r}(u))\,\mathrm{d}u
    \right)\text{。}
 \end{aligned}
 \label{eq:vision-nerf-continuous}
\end{equation}
$T(s)$称为\term{透射率}（transmittance），表示光从相机出发到达距离$s$之前没有被截获的比例；近处密度越大，远处颜色的可见贡献越小。$\sigma$的单位是长度的倒数，使指数无量纲。

数值计算时沿射线取$n$个有序样本，假定每段内密度和颜色不变。设第$i$段长为$\delta_i$，得到
\begin{equation}
 \begin{aligned}
 \alpha_i&=1-\exp(-\sigma_i\delta_i),&
 T_i&=\prod_{j<i}(1-\alpha_j),\\
 w_i&=T_i\alpha_i,&
 \widehat{\bm{C}}&=\sum_{i=1}^{n}w_i\bm{c}_i
       +T_{n+1}\bm{c}_{\mathrm{bg}}\text{。}
 \end{aligned}
 \label{eq:vision-nerf-discrete}
\end{equation}
$\alpha_i$是光在第$i$段被截获的概率，空乘积取1，$\bm{c}_{\mathrm{bg}}$是背景色。权重$w_i$由“前面各段都没有截获”和“本段截获”两部分相乘得到，所以远处的颜色不能只按自身的$\alpha_i$相加。各$w_i$与背景权重$T_{n+1}$之和为1。

用一个两段的例子手算。取$\alpha_1=0.5,\alpha_2=0.75$，颜色分别为红$(1,0,0)$和蓝$(0,0,1)$，背景为白$(1,1,1)$。第一段权重为0.5；第二段只能使用剩下的0.5，权重为$0.5\times0.75=0.375$；背景权重为$0.5\times0.25=0.125$。合成颜色为$(0.625,0.125,0.5)$。若忘记透射率，直接用0.5和0.75加权，就会重复计算被前景遮住的贡献，权重和也超过1。图\ref{fig:vision-ray-compositing}画出了这个顺序和权重。

\begin{figure}[htbp]
 \centering
 % Two intervals, alpha .5 and .75; background transmittance .125.
\begin{tikzpicture}[x=1mm,y=1mm,font=\figurefont\small]
 \draw[BookMuted,line width=.9pt,->] (5,0)--(100,0) node[right] {$s$};
 \draw[BookInk,line width=.9pt] (5,-3)--(0,-6)--(0,6)--(5,3)--cycle;
 \node[anchor=south] at (4,9) {相机};
 \fill[BookBlue!12] (24,-7) rectangle (41,7);
 \draw[BookBlue,line width=.9pt] (24,-7) rectangle (41,7);
 \fill[BookTeal!15] (58,-7) rectangle (75,7);
 \draw[BookTeal,line width=.9pt] (58,-7) rectangle (75,7);
 \draw[BookMuted,line width=1pt] (91,-7)--(91,7);
 \node[anchor=south] at (32.5,8) {$\alpha_1=0.5$};
 \node[anchor=south] at (66.5,8) {$\alpha_2=0.75$};
 \node[anchor=north] at (32.5,-9) {$T_1=1$};
 \node[anchor=north] at (66.5,-9) {$T_2=0.5$};
 \node[anchor=south] at (91,8) {背景};
 \node[anchor=east] at (7,-31) {$w$};
 \fill[BookBlue!20] (10,-35) rectangle (50,-27);
 \fill[BookTeal!20] (50,-35) rectangle (80,-27);
 \fill[BookPaper] (80,-35) rectangle (90,-27);
 \draw[BookInk,line width=.7pt] (10,-35) rectangle (90,-27);
 \draw[BookInk,line width=.7pt] (50,-35)--(50,-27);
 \draw[BookInk,line width=.7pt] (80,-35)--(80,-27);
 \node at (30,-31) {$0.5$};
 \node at (65,-31) {$0.375$};
 \node[anchor=north] at (30,-37) {第一段};
 \node[anchor=north] at (65,-37) {第二段};
 \node[anchor=north] at (85,-37) {$0.125$};
\end{tikzpicture}

 \caption{两段射线的体渲染教学例。近处段先取得0.5的贡献权重，后段只能使用剩余0.5中的0.75，背景使用最终透射率0.125。下方长度条表示实际权重，三段和为1。颜色值人为设定，用于复算，图中形状不表示真实物体测量。}
 \label{fig:vision-ray-compositing}
\end{figure}

网络本身的结构如下。多层感知机直接读取三维坐标时倾向于学到低频函数，难以表示细纹理，因此先对输入作多频正余弦编码：每个坐标分量$p$映射为$(\sin2^0\pi p,\cos2^0\pi p,\ldots,\sin2^{L-1}\pi p,\cos2^{L-1}\pi p)$，位置取$L=10$得到60维，方向取$L=4$得到24维。这种位置编码（positional encoding）使网络能够拟合高频细节。编码前，位置的三个坐标先归一化到$[-1,1]$，方向取单位向量，各分量本身也在$[-1,1]$内。网络的逐层结构见表\ref{tab:vision-nerf-structure}：位置编码经8层宽256的全连接层（ReLU激活），按DeepSDF的做法，第5层的输出再与输入编码拼接后送入第6层；之后一层同时输出经ReLU保证非负的密度$\sigma$和一个不加激活的256维特征；该特征与方向编码拼接，经一层宽128的全连接层（ReLU）和一层sigmoid输出层得到RGB颜色。密度不读取方向，保证几何在各视角下一致。

\begin{table}[htbp]
 \centering
 \small
 \begin{tabularx}{\linewidth}{@{}p{16mm}Xp{15mm}p{9mm}p{6mm}p{15mm}@{}}
  \toprule
  阶段或层组 & 层类型 & 卷积核与步长 & 输出通道 & 重复次数 & 输出尺寸\\
  \midrule
  位置编码 & $\gamma(\bm{x})$，$L=10$ & — & 60 & 1 & 每个采样点\\
  主干1至5 & 全连接，ReLU & — & 256 & 5 & 每个采样点\\
  跳连拼接 & 第5层输出与$\gamma(\bm{x})$拼接 & — & 316 & 1 & 每个采样点\\
  主干6至8 & 全连接，ReLU & — & 256 & 3 & 每个采样点\\
  密度与特征 & 全连接；$\sigma$经ReLU，特征无激活 & — & 1与256 & 1 & 每个采样点\\
  方向分支 & 特征与$\gamma(\bm{d})$（$L=4$，24维）拼接，全连接，ReLU & — & 128 & 1 & 每个采样点\\
  颜色输出 & 全连接，sigmoid & — & 3 & 1 & 每个采样点\\
  \bottomrule
 \end{tabularx}
 \caption{NeRF的多层感知机，按原论文正文与附录图7整理。网络逐点计算，没有卷积，故卷积核一列为空、输出尺寸为单个采样点。原论文同时优化“粗”“细”两个网络，参数各自独立；论文只给出一种网络结构，没有区分二者。论文报告单个场景的网络权重约5\,MB。}
 \label{tab:vision-nerf-structure}
\end{table}

逐点查询网络很昂贵，而射线上大部分点位于空处或被遮挡。NeRF因此用粗细两级采样：粗网络在射线上分层随机取64个点，计算权重$w_i$，把它们归一化成沿射线的分段分布；再按这个分布多采128个点，细网络在全部192个点上重新计算颜色。粗细两个网络的颜色都与真实像素比较。训练时每次从所有训练照片中随机抽4096条像素射线，损失是粗、细渲染颜色与真实颜色的平方误差之和，用Adam优化，学习率从$5\times10^{-4}$指数衰减到$5\times10^{-5}$，其余参数取默认值$\beta_1=0.9,\beta_2=0.999,\epsilon=10^{-7}$；原论文一个场景训练10万到30万次迭代，在一张V100上需要1到2天。

渲染新视角时，为新相机的每个像素构造射线，重复上述采样与积分。一幅图有$p$条射线、每条查询$n$个点，计算量随$pn$增长。渲染时每条射线经粗网络查询64个点、经细网络查询192个点，共256次；原论文的合成场景每幅图有64万条射线，真实场景有76.2万条，每幅图需要1.5亿到2亿次网络查询，在一张V100上约30秒一帧，远不能实时渲染。图像损失本身不唯一确定几何：未拍到的表面、镜面反射和错误的相机都可能让新视角失败，而训练视图仍然拟合得很好。

\subsubsection{3D Gaussian Splatting}
\label{sec:vision-generation-gaussians}

NeRF的瓶颈在于每个像素都要沿射线查询上百次网络，渲染速度远不能实时，训练也要一两天。大部分查询落在空处，却都要付出完整的网络计算。3D Gaussian Splatting的动机是用显式的几何单元只在有物体的地方表示场景，并用光栅化代替射线采样，从而实现实时渲染和更快的训练\scite{vision-gaussian2023}。

主要想法是：用一组三维高斯椭球表示场景，每个高斯有位置、形状、不透明度和随方向变化的颜色；渲染时把每个高斯投影成图像上的二维椭圆，按深度从前到后做透明度合成。这个合成公式与NeRF的离散体渲染具有相同的结构，但只在高斯覆盖的像素上计算，不需要沿射线查询网络。

每个高斯的参数包括：中心$\bm{\mu}\in\mathbb{R}^3$；协方差$\bm{\Sigma}=\bm{R}_{g}\bm{S}\bm{S}^{\top}\bm{R}_{g}^{\top}$，其中旋转$\bm{R}_{g}$由单位四元数（4个数）表示、$\bm{S}$是由3个正尺度组成的对角阵，这种分解保证协方差始终半正定；不透明度$o\in(0,1)$，经sigmoid得到；颜色用3阶\term{球谐函数}（spherical harmonics）系数表示，即在方向球面上用16个基函数的组合表示每个颜色通道随方向的变化，共48个系数。一个高斯共有$3+4+3+1+48=59$个参数，一个场景通常有一百万到五百万个高斯。论文没有把59这个总数写出来，它是按上述参数化逐项相加得到的；场景规模“一百万到五百万个高斯”是论文对所测全部场景的报告。各参数的维数、约束方式与初始化见表\ref{tab:vision-gaussian-parameters}。

\begin{table}[htbp]
 \centering
 \small
 \begin{tabularx}{\linewidth}{@{}p{20mm}p{9mm}XX@{}}
  \toprule
  参数 & 维数 & 约束方式 & 初始化\\
  \midrule
  中心$\bm{\mu}$ & 3 & 无约束；学习率按指数衰减 & 运动恢复结构的稀疏点\\
  旋转$\bm{q}$ & 4 & 单位四元数，使用前归一化 & 论文未给出\\
  尺度$\bm{s}$ & 3 & 经指数函数得到正尺度 & 各向同性，取到最近三个点的平均距离\\
  不透明度$o$ & 1 & 经sigmoid落在$[0,1)$ & 论文未给出\\
  球谐系数 & 48 & 3阶，每个颜色通道16个；先只优化0阶，每1000次迭代开放一阶 & 论文未给出\\
  \bottomrule
 \end{tabularx}
 \caption{3D Gaussian Splatting中单个高斯的参数，取自原论文。59个参数在每个高斯上独立优化，场景中高斯的数目由自适应密度控制在训练中改变。}
 \label{tab:vision-gaussian-parameters}
\end{table}

渲染时，先把三维高斯变换到相机坐标，再在中心附近把透视投影线性化，得到二维协方差
\begin{equation}
 \bm{\Sigma}_{2}
 =\bm{J}\bm{R}_{c}\bm{\Sigma}
    \bm{R}_{c}^{\top}\bm{J}^{\top}\text{，}
 \label{eq:vision-gaussian-projection}
\end{equation}
其中$\bm{R}_{c}$把世界方向转到相机坐标，$\bm{J}\in\mathbb{R}^{2\times3}$是该点处投影的Jacobian矩阵。透视投影不是线性的，所以这只是局部近似。设投影中心为$\bm{\mu}_{2,i}$，第$i$个高斯在像素$\bm{u}$处的有效不透明度为
\begin{equation}
 a_i(\bm{u})=o_i\exp\!\left[
 -\frac12(\bm{u}-\bm{\mu}_{2,i})^\top
 \bm{\Sigma}_{2,i}^{-1}(\bm{u}-\bm{\mu}_{2,i})
 \right]\text{。}
 \label{eq:vision-gaussian-alpha}
\end{equation}
把覆盖该像素的高斯按深度从近到远排列，颜色为$\sum_i a_i\bm{c}_i\prod_{j<i}(1-a_j)$，必要时加上剩余背景。与式\eqref{eq:vision-nerf-discrete}对照，$a_i$对应$\alpha_i$，$\prod_{j<i}(1-a_j)$对应透射率$T_i$；区别在于$a_i$来自投影椭圆和学到的不透明度，而不是对密度的分段积分。

为了在GPU上高效完成，原方法把屏幕分成$16\times16$像素的图块，每个高斯按它覆盖的图块复制若干份，用“图块编号加深度”作为键做一次全局基数排序，然后每个图块并行地按深度合成。排序按高斯中心的深度进行，同一像素内的次序只是近似，不做精确的表面求交。

初始化来自运动恢复结构给出的稀疏点：每个点成为一个高斯，初始为各向同性，尺度取到最近三个点的平均距离。损失为$L=(1-\lambda)L_1+\lambda L_{\text{D-SSIM}}$，$\lambda=0.2$，其中D-SSIM是由结构相似性构造的差异项，公开实现取1减SSIM。为了稳定，优化先在缩小为四分之一的图像分辨率上进行，第250次和第500次迭代后各放大一倍。优化同时更新所有高斯的位置、尺度、旋转、不透明度和球谐系数；球谐阶数从0开始，每1000次迭代增加一阶，直到3阶。

稀疏点不足以表达所有细节，因此需要\term{自适应密度控制}（adaptive density control），在训练中改变高斯的数目。优化热身之后每100次迭代检查一次：视图空间位置梯度的平均幅值超过阈值0.0002的高斯，说明该处重建不足。若高斯较小，就复制一份并沿梯度方向移动，用于覆盖缺失的结构；若高斯较大，就分裂成两个，尺度各除以1.6，新位置从原高斯中采样，用于细化过大的区域。不透明度低于阈值$\epsilon_\alpha$的高斯被删除（论文未给出$\epsilon_\alpha$的数值），世界空间中过大或在视图中覆盖范围过大的高斯也定期删除；每3000次迭代把所有不透明度重置为接近零，让多余的高斯在后续优化中自然淘汰。例如玩具狗的耳朵边缘若被一个大高斯覆盖，耳朵与背景的颜色会被混在一起；分裂后两个较小的高斯可以分别拟合耳朵和背景，但这仍需后续优化，分裂本身不代表测得了新的表面点。原论文训练3万次迭代，在Mip-NeRF360数据集的场景上约需40分钟。

渲染新视角只需重新投影、排序和合成，原论文在这些场景上以超过100帧每秒的速度渲染，而对照的Mip-NeRF360每帧需要数秒以上。代价是存储：一个场景模型约数百兆字节，而NeRF类模型只需几兆字节。高斯集合也不是带拓扑的三角网格，没有物理材质，渲染逼真不说明它可以直接用于碰撞检测或精确编辑。

两种场景表示的比较见表\ref{tab:vision-generation-scene-comparison}。固定任务为“从绕玩具狗拍摄的照片渲染新视角”，评价口径为保留视图上的画质、训练与渲染时间以及存储。3D Gaussian Splatting相对NeRF改变了表示（隐式网络改为显式高斯）和渲染方式（射线采样改为投影排序），换来的是实时渲染和以分钟计的训练，付出的是数百倍的存储和对初始点云的依赖；两者都无法为没有拍到的表面提供观测支持。

\begin{table}[htbp]
 \centering
 \small
 \begin{tabularx}{\linewidth}{@{}l X X@{}}
  \toprule
   & NeRF & 3D Gaussian Splatting\\
  \midrule
  场景表示 & 多层感知机隐式表示的密度与颜色场 & 百万量级的显式三维高斯\\
  渲染方式 & 每条射线查询网络两百余次 & 投影、按图块排序、透明度合成\\
  初始化 & 随机初始化网络 & 运动恢复结构的稀疏点\\
  训练时间（原论文设置） & 单场景1到2天 & 单场景约40分钟\\
  渲染速度 & 逐射线查询，远低于实时 & 每秒百帧以上\\
  存储 & 几兆字节量级的网络权重 & 数百兆字节的高斯参数\\
  \bottomrule
 \end{tabularx}
 \caption{两种场景表示在同一新视角合成任务下的比较。时间与存储取自各自原论文报告的设置，硬件与数据集不同，只表示数量级差别。两者都需逐场景优化，要求静态场景与准确相机，且不能为未观测表面提供证据。}
 \label{tab:vision-generation-scene-comparison}
\end{table}

\subsection{新视角的渲染与核验}
\label{sec:vision-generation-rendering-check}

渲染一个新视角时，给出目标相机的内参和外参，按上面的方法得到图像。核验的关键是使用在拟合前就留出的测试视图：在这些相机下渲染，与真实照片比较PSNR、SSIM或LPIPS（见\ref{chap:vision-image-restoration}章），同时检查轮廓、细节和反光。测试视图的像素不能用于调节场景参数；若相机标定使用了测试图像，也应在协议中说明。

单个测试视图还可能漏掉不连续现象。沿一条连续的相机路径渲染视频，可以观察遮挡变化时是否出现漂浮的斑点、外观是否突然跳变。路径上没有真实照片的位置只能做一致性观察，无法计算误差。几何精度还需要点云、深度等独立测量，逼真连贯的渲染动画不能直接证明表面准确。对玩具狗的背面，若拍摄没有覆盖，两种表示都可能把错误拟合延伸过去。若用途只是展示某些可见视角，应注明适用的视角范围；若需要完整的三维资产，就需要补拍背面并核验几何。动态场景的时间变化由\ref{chap:vision-multimodal-generation}章扩展。

\section{生成结果的评价}
\label{sec:vision-generation-evaluation}

前四节分别给出了各自任务的核验方法，本节把它们放在同一组维度下比较。四类任务的参照不同：创建主要对照条件，转换和编辑还要对照输入，新视角要对照真实观测。相应地，评价分为三个维度：条件是否被遵循，输入中应保持的内容是否被保持，以及结果有多少受观测支持、付出了多少成本。\ref{sec:vision-generation-challenges}节定义的FID、CLIP分数、人工偏好、保留区变化和PSNR等指标分别落在这三个维度之下，本节讨论怎样组合使用它们。比较模型时还要固定提示、参考、分辨率、候选数量和修订预算，否则无法判断差异来自模型本身还是选择过程。

\subsection{条件遵循的核验}
\label{sec:vision-generation-condition-check}

回到玩具狗与杯子的要求记录，逐项检查存在、数量、颜色归属和左右关系；对主体适配，还要与参考照片比较身份。可以用检测器统计对象和位置，用区域特征比较主体，用文字读取器检查杯子上的字符。DreamBooth的评价是一个例子：它用DINO特征相似度衡量主体保真、用CLIP图文相似度衡量提示遵循，两个指标同时报告，防止“复制参考照片”这种退化在单一指标上得高分。这些检查工具本身也会漏检或误读，自动分数应注明工具的输入和输出，并抽样人工复核，不能把工具的错误算到生成器头上。

整体图文相似度适合衡量粗粒度相关性，不能单独证明复杂条件全部成立。汇总时除平均分外，应保留关键分项和失败率，例如对象数量正确率、属性绑定正确率。多个属性中一个绑定错误，即使只影响一小块区域，也会使整个请求失败。对使用视觉指令的编辑，还要检查结果是否随指令改变，而不只是检查目标对象是否出现。

\subsection{输入保持的核验}
\label{sec:vision-generation-change-check}

转换和编辑都需要同时检查“该改的改了”和“该留的留了”。原样返回原图，可以使式\eqref{eq:vision-edit-preservation}的保留误差为零，却没有执行请求；把杯子换成花瓶的同时改变了狗的斑纹，完成了目标却破坏了保留条件。两类错误应分别记录。多轮编辑还应分别报告当前轮修改的完成率、历次已接受修改的保持情况，以及相对最初主体的漂移。

保持的含义随任务而变。照片转成绘画后，像素距离大是预期的现象，应该比较轮廓、对象数量和身份，而不是逐像素误差；LPIPS这类感知距离也要结合允许的风格变化来解释。新视角中，同一物体因投影而在画面中移动，也不是编辑漂移；评价必须与目标相机下的真实照片比较，而不是与原相机的图像逐像素相减。

\subsection{观测支持与制作成本的核验}
\label{sec:vision-generation-costs}

观测支持回答“结果中的哪些内容有证据”。新视角合成的输出应标明哪些区域被训练视图覆盖；DreamBooth生成的背面若没有参考照片，就只是类别先验的推测；修补区域中的花瓶是按请求创作的内容，而不是被遮挡物体的恢复。三者都可以产生逼真图像，但真实性的依据不同，交付时应分别说明。

成本则回答“得到这一幅结果花了多少”。一项可交付成果的成本包括条件制作（画姿态图、写掩码）、模型适配（ControlNet支路训练、DreamBooth微调）、视图采集与标定、场景拟合、候选生成与筛选以及修订轮次。一次渲染只需几毫秒，不表示完成整个场景也只需一次前向；从一百幅候选中选出一幅，也不表示单次生成就稳定满足要求。报告时应把这些成本与结果一起列出，使不同系统在相同预算下比较。

\section{本章小结}
\label{sec:vision-generation-summary}

开篇的四个请求，现在可以分别回答。它们共同面对四个挑战：一个条件对应许多合格图像，多个对象的属性容易串换，局部修改与内容保持互相牵制，有限的照片难以保证三维一致；评价时则需要分别用FID、CLIP分数、人工偏好、保留区变化和PSNR等指标回答不同的问题。在雪地里画这只玩具狗，属于新图像的创建：文字经LDM或unCLIP约束语义，ControlNet通过冻结主干、复制支路和零卷积接入逐像素对齐的空间图，DreamBooth用少量照片和类别先验保持把具体主体写进模型参数，多参考生成则需要显式规定每幅参考的角色。采样中的引导以多样性和计算为代价提高条件贴合度，生成后的重排只能在已有候选中选择，一轮反思修订可以纠正某些空间关系，但需要预算上限和独立验收。转成水彩画属于已有图像的转换：pix2pix依赖配对目标，CycleGAN用循环一致性替代配对却不能唯一确定语义，Johnson感知损失网络用固定特征网络把风格转换变为一次前向计算。把杯子换成花瓶属于指定内容的编辑：泊松图像编辑在梯度域中把现成的花瓶无缝贴入，Criminisi样例修补按置信度与等照度线强度决定填充顺序、用图中已有的块补洞，PatchMatch用随机初始化、传播和随机搜索把最近邻块的搜索加速到可交互，这三种经典方法让保留区严格不变，却只能重组已有内容；GLIDE修补模型靠掩码与可见区域约束范围，InstructPix2Pix用合成的编辑三元组和双条件引导实现只需一句指令的编辑，多轮编辑还要在模型之外维护版本状态。给出正后方视图属于新视角合成：多视图几何先给出相机与稀疏点，光场渲染需要密集拍摄才能直接插值光线；NeRF用隐式场与体渲染、3D Gaussian Splatting用显式高斯与投影合成拟合同一组观测，后者以存储换来实时渲染。

评价贯穿这四类任务：条件遵循、输入保持、观测支持与成本应分别核验，不能用一个自然度分数或全局相似度代替。仍有一些问题本章没有解决。多主体和多参考的角色绑定还缺少经过验证的通用模型；精确数量、文字和复杂空间关系仍是文字条件生成的常见失败；多轮编辑的主体漂移只能靠外部回退限制，不能从模型内部消除；缺少视角的观测场景无法由本章的拟合方法补全，需要\ref{chap:vision-multimodal-generation}章引入生成先验，那时又必须重新区分哪些内容有观测依据。视频生成还会增加时间连续性和音画同步的约束，本章关于条件、保持与观测证据的区分在那里仍然适用。

\paragraph{扩展阅读}
创作条件方向，SDXL在LDM的基础上扩大去噪网络，改用两个文字编码器，并把训练图像的原始尺寸和裁剪位置也作为条件输入，再用一个细化模型在同一潜空间中补充细节，可用来观察本章的LDM结构在规模扩大后还需要补充哪些条件\scite{vision-sdxl2024}T2I-Adapter同样向冻结的文生图模型接入草图、深度或姿态等空间图，但只训练一个轻量的适配网络，不像ControlNet那样复制编码器支路，多个适配器的输出还可以相加组合\scite{vision-t2i-adapter2024}IP-Adapter为参考图像单独设置一组与文字交叉注意力并行的交叉注意力，使参考图像不经微调即可作为条件；Textual Inversion则冻结整个模型，只从少量照片中学习一个新词的嵌入向量。二者与DreamBooth一样要把具体主体或风格带入生成，但分别把信息放在推理时的条件和词嵌入中，而不是写进模型参数，保真程度和成本也随之不同\scite{vision-ip-adapter2023,vision-textual-inversion2023}

编辑方向，SDEdit先把用户的笔画草图或待改图像加噪到某个中间时刻，再用扩散模型去噪回来；加噪越多，结果越自然，与输入的对应却越弱。它不需要配对数据或额外训练，可与本章pix2pix和CycleGAN的学习式转换对照\scite{vision-sdedit2022}本章的InstructPix2Pix只在两幅生成图像之间使用Prompt-to-Prompt；Null-text Inversion先用确定性采样把真实照片反演为一条去噪轨迹，再逐步优化无分类器引导中空文本条件的嵌入，使加了引导的采样仍能准确重建原图，从而把同样的注意力控制用于真实照片的编辑\scite{vision-null-text-inversion2023}新视角合成方向，Mip-NeRF把每个像素对应的射线换成锥体，用积分位置编码描述锥体截面覆盖的区域，缓解NeRF在拍摄距离和分辨率变化时的混叠；Instant-NGP用多分辨率哈希表存储可学习的特征，把NeRF类模型的训练从一两天缩短到秒级至分钟级，从隐式表示一侧回应了3D Gaussian Splatting所针对的速度问题\scite{vision-mip-nerf2021,vision-instant-ngp2022}

评价方向，Kynkäänniemi等人发现，只要让生成图像在ImageNet类别上的分布与真实图像对齐，即使画质没有改善，FID也会明显下降，说明FID在相当程度上反映的是Inception特征中的类别信息\scite{vision-fid-imagenet-classes2023}Jayasumana等人认为Inception特征和正态分布假设都不适合评价现代文生图模型，改用CLIP特征和不假设分布形式的最大均值差异，提出CMMD\scite{vision-cmmd2024}GenEval用目标检测器核对生成图像中的物体是否出现，以及数量、颜色和相对位置是否符合提示，把本章所说的属性绑定和计数错误变成可以自动统计的检查项\scite{vision-geneval2023}
